% 第4章 理想流体动力学 —— 依据原书扫描页手工精修版
\chapter{理想流体动力学}

本章详细介绍忽略流体粘性的理想流体运动的理论分析和计算方法以及它们的主要性质.

\section*{4.1  理想流体运动的基本方程和初边值条件}\addcontentsline{toc}{section}{4.1 理想流体运动的基本方程和初边值条件}

理想流体是真实流体的一种近似模型，这种近似使流体动力学问题大为简化，到19世纪末，理想流体动力学已经形成完整的理论体系.历史上，20世纪以前，所谓的水动力学(Hydrodynamics)实际上是理想流体动力学.然而，当时工程界对理想流体动力学理论是否符合实际常常提出质疑.理想流体动力学的最主要缺陷是不能预测在流体中运动物体的阻力.理想流体动力学从纯理论的学科成为有广泛工程应用的科学分支始于20世纪初.德国流体力学家普朗特发现了绕流物体表面的边界层现象，这一现象揭示了“稀流体”(粘性较小的流体)绕流线形物体运动时，粘性只在物体表面很薄一层里才对流动有显著的影响，忽略边界层厚度和流体粘性而用理想流体动力学理论来预测流线形物体表面上的压强分布和实际情况符合得很好.这一发现不仅说明了理想流体动力学理论的适用范围，而且在飞机工业的发展中起了推动作用，最主要的早期贡献是用理想流体动力学理论正确解释了飞机机翼上升力产生的原因.迄今为止，对于没有流动分离的绕流现象，理想流体动力学是关于预测流场和物面压强分布相当好的理论和方法.至于超出理想流体近似范围以外的流动现象，需要应用粘性流体力学知识来解释和研究.本章先介绍理想流体动力学，第8章介绍粘性流体力学.

理想流体包括理想不可压缩流体(液体)和理想可压缩气体.

\subsection*{1. 理想流体运动的基本方程——欧拉方程}

理想流体忽略粘性，它的应力状态是各向同性的，即
\begin{equation*}
T_{ij}=-p\,\delta_{ij}
\tag{4.1}
\end{equation*}
将它代入第3章导出的流体动力学守恒方程，得到向量形式的理想流体动力学方程组如下：

连续性方程
\begin{equation*}
\frac{\partial\rho}{\partial t}+\nabla\cdot\rho\boldsymbol{U}=0
\tag{4.2a}
\end{equation*}
运动方程
\begin{equation*}
\frac{\partial\boldsymbol{U}}{\partial t}+\boldsymbol{U}\cdot\nabla\boldsymbol{U}=-\frac{1}{\rho}\,\nabla p+\boldsymbol{f}
\tag{4.2b}
\end{equation*}
能量方程
\begin{equation*}
\frac{\partial}{\partial t}\Big(e+\frac{|\boldsymbol{U}|^{2}}{2}\Big)+\boldsymbol{U}\cdot\nabla\Big(e+\frac{|\boldsymbol{U}|^{2}}{2}\Big)=-\frac{1}{\rho}\,\nabla\cdot(p\boldsymbol{U})+\boldsymbol{f}\cdot\boldsymbol{U}+\dot{q}
\tag{4.2c}
\end{equation*}
在直角坐标系中上述方程表示为

连续性方程
\begin{equation*}
\frac{\partial\rho}{\partial t}+\frac{\partial\rho u}{\partial x}+\frac{\partial\rho v}{\partial y}+\frac{\partial\rho w}{\partial z}=0
\tag{4.3a}
\end{equation*}
\begin{equation*}
\begin{cases}
\dfrac{\partial u}{\partial t}+u\dfrac{\partial u}{\partial x}+v\dfrac{\partial u}{\partial y}+w\dfrac{\partial u}{\partial z}=-\dfrac{1}{\rho}\,\dfrac{\partial p}{\partial x}+f_x\\[8pt]
\dfrac{\partial v}{\partial t}+u\dfrac{\partial v}{\partial x}+v\dfrac{\partial v}{\partial y}+w\dfrac{\partial v}{\partial z}=-\dfrac{1}{\rho}\,\dfrac{\partial p}{\partial y}+f_y\\[8pt]
\dfrac{\partial w}{\partial t}+u\dfrac{\partial w}{\partial x}+v\dfrac{\partial w}{\partial y}+w\dfrac{\partial w}{\partial z}=-\dfrac{1}{\rho}\,\dfrac{\partial p}{\partial z}+f_z
\end{cases}
\tag{4.3b}
\end{equation*}
能量方程
\begin{align*}
\frac{\partial}{\partial t}\Big(e+\frac{|\boldsymbol{U}|^{2}}{2}\Big)+u\,\frac{\partial}{\partial x}\Big(e+\frac{|\boldsymbol{U}|^{2}}{2}\Big)+v\,\frac{\partial}{\partial y}\Big(e+\frac{|\boldsymbol{U}|^{2}}{2}\Big)+w\,\frac{\partial}{\partial z}\Big(e+\frac{|\boldsymbol{U}|^{2}}{2}\Big)\\
=-\frac{1}{\rho}\bigg[\frac{\partial(pu)}{\partial x}+\frac{\partial(pv)}{\partial y}+\frac{\partial(pw)}{\partial z}\bigg]+u f_x+v f_y+w f_z+\dot{q}
\tag{4.3c}
\end{align*}
理想流体运动的控制方程又称欧拉方程.

对于均质不可压缩流体，它的密度$\rho=\mathrm{const.}$，因此式(4.2a)，式(4.2b)可简化为
\begin{equation*}
\nabla\cdot\boldsymbol{U}=0
\tag{4.4a}
\end{equation*}
\begin{equation*}
\frac{\partial\boldsymbol{U}}{\partial t}+\boldsymbol{U}\cdot\nabla\boldsymbol{U}=-\frac{1}{\rho}\,\nabla p+\boldsymbol{f}
\tag{4.4b}
\end{equation*}
由于不可压缩流体微团的体积膨胀率等于零，因此微团上的体积膨胀功也为零，从而微团的内能增长率等于外界输入的热量，即$\mathrm{d}e/\mathrm{d}t=\dot{q}$，于是能量方程(4.3c)简化为
\begin{equation*}
\frac{\partial}{\partial t}\Big(\frac{|\boldsymbol{U}|^{2}}{2}\Big)+\boldsymbol{U}\cdot\nabla\Big(\frac{|\boldsymbol{U}|^{2}}{2}\Big)=-\frac{1}{\rho}\,\nabla\cdot(p\boldsymbol{U})+\boldsymbol{f}\cdot\boldsymbol{U}
\tag{4.4c}
\end{equation*}
上式也可以直接由运动方程式(4.4b)点乘速度向量$\boldsymbol{U}$后导出，因此不可压缩理想流体的能量方程也就是运动方程的能量积分，它不再是独立的方程.事实上，不可压缩流体的密度不随压强变化，是给定的物性参数，所以方程组(4.4a)，(4.4b)已组成封闭方程组.

对于完全气体，有以下的热力学状态方程：
\begin{gather*}
p=R\rho T\tag{4.5a}\\[2pt]
e=c_v T\tag{4.5b}
\end{gather*}
其中$R$为气体常数，由下式算出：
\begin{equation*}
R=\frac{8314}{\mu}
\tag{4.6}
\end{equation*}
$R$的单位是$\mathrm{N\cdot m\cdot kg^{-1}\cdot K^{-1}}$，$\mu$是气体的相对分子质量.例如，空气的相对分子质量是29，$R=\dfrac{8314}{29}=287$.式(4.5b)中，$c_v$是气体的比定容热容，对于完全气体来说
\begin{gather*}
c_v=\frac{R}{\gamma-1}\tag{4.7}\\[2pt]
\gamma=\frac{c_p}{c_v}\tag{4.8}
\end{gather*}
式中，$c_p$是气体的比定压热容，$\gamma$称作比热容比或绝热指数.方程组(4.2a)、(4.2b)，(4.2c)和(4.5a)，(4.5b)联立构成完全气体运动的基本方程组.

\subsection*{2. 初边值条件}

(1)初始条件

对于非定常流动，应给出流场的初始状态，即$t=0$时的速度场、压强场和密度场：
\begin{equation*}
\boldsymbol{U}(\boldsymbol{x},0)=\boldsymbol{U}(\boldsymbol{x})\,,\qquad p(\boldsymbol{x},0)=p(\boldsymbol{x})\,,\qquad \rho(\boldsymbol{x},0)=\rho(\boldsymbol{x})
\tag{4.9}
\end{equation*}
对于定常流动，任何时刻流动状态都相同，基本方程中局部导数项等于零，这时就不需要给出初始条件.

(2)边界条件

\textcircled{1} 固体壁面的不可穿透条件

理想流体忽略粘性，在固体壁面上不存在剪应力，因此理想流体质点在壁面上可以任意滑动.固体壁面的约束条件是：垂直于壁面的法向速度连续，即
\begin{equation*}
(\boldsymbol{U}\cdot\boldsymbol{n})_\Sigma=(\boldsymbol{U}_b\cdot\boldsymbol{n})_\Sigma
\tag{4.10}
\end{equation*}
$\boldsymbol{U}_b$为固体壁面上任一点的移动速度，$\boldsymbol{U}$为同一点流体质点速度，$\boldsymbol{n}$为该点固壁面的法向量.当给出壁面$\Sigma$的解析式时：
\begin{equation*}
\Sigma:F(x,y,z,t)=0
\tag{4.11}
\end{equation*}
壁面条件(4.10)可以由壁面方程表示为(直角坐标)
\begin{equation*}
\frac{\partial F}{\partial t}+u\,\frac{\partial F}{\partial x}+v\,\frac{\partial F}{\partial y}+w\,\frac{\partial F}{\partial z}=0
\tag{4.12}
\end{equation*}
式中$u,v,w$是壁面上流体质点速度分量.式(4.12)可证明如下，对壁面函数式(4.11)求全导数，得
\[
\frac{\mathrm{d}F}{\mathrm{d}t}=\frac{\partial F}{\partial t}+\frac{\partial F}{\partial x}\,\frac{\mathrm{d}x}{\mathrm{d}t}+\frac{\partial F}{\partial y}\,\frac{\mathrm{d}y}{\mathrm{d}t}+\frac{\partial F}{\partial z}\,\frac{\mathrm{d}z}{\mathrm{d}t}=0
\]
式中$\mathrm{d}x/\mathrm{d}t,\mathrm{d}y/\mathrm{d}t,\mathrm{d}z/\mathrm{d}t$是壁面上任意点$(x,y,z)$的速度分量；壁面上任一点的法向量可用壁面方程表示如下：
\begin{equation*}
\boldsymbol{n}=\frac{\nabla F}{|\nabla F|}=\frac{\dfrac{\partial F}{\partial x}\boldsymbol{e}_x+\dfrac{\partial F}{\partial y}\boldsymbol{e}_y+\dfrac{\partial F}{\partial z}\boldsymbol{e}_z}{\sqrt{\Big(\dfrac{\partial F}{\partial x}\Big)^{2}+\Big(\dfrac{\partial F}{\partial y}\Big)^{2}+\Big(\dfrac{\partial F}{\partial z}\Big)^{2}}}
\tag{4.13}
\end{equation*}
于是壁面速度的法向分量为
\[
\boldsymbol{U}_B\cdot\boldsymbol{n}=\frac{\dfrac{\partial F}{\partial x}\,\dfrac{\mathrm{d}x}{\mathrm{d}t}+\dfrac{\partial F}{\partial y}\,\dfrac{\mathrm{d}y}{\mathrm{d}t}+\dfrac{\partial F}{\partial z}\,\dfrac{\mathrm{d}z}{\mathrm{d}t}}{\sqrt{\Big(\dfrac{\partial F}{\partial x}\Big)^{2}+\Big(\dfrac{\partial F}{\partial y}\Big)^{2}+\Big(\dfrac{\partial F}{\partial z}\Big)^{2}}}=\frac{-\dfrac{\partial F}{\partial t}}{\sqrt{\Big(\dfrac{\partial F}{\partial x}\Big)^{2}+\Big(\dfrac{\partial F}{\partial y}\Big)^{2}+\Big(\dfrac{\partial F}{\partial z}\Big)^{2}}}
\]
位于壁面上的流体质点速度法向分量为
\[
\boldsymbol{U}\cdot\boldsymbol{n}=-\frac{u\,\dfrac{\partial F}{\partial x}+v\,\dfrac{\partial F}{\partial y}+w\,\dfrac{\partial F}{\partial z}}{\sqrt{\Big(\dfrac{\partial F}{\partial x}\Big)^{2}+\Big(\dfrac{\partial F}{\partial y}\Big)^{2}+\Big(\dfrac{\partial F}{\partial z}\Big)^{2}}}
\]
由于$(\boldsymbol{U}\cdot\boldsymbol{n})_\Sigma=(\boldsymbol{U}_B\cdot\boldsymbol{n})_\Sigma$，可得式(4.12).

第2章中曾经指出，运动学的公式和坐标系的选择有关.下面给出几种常用的固壁上流体运动学条件的具体表达式.

(i) 在固定坐标系中静止固壁的边界条件(见图4.1)

静止固壁的物面方程和时间无关，即有$\dfrac{\partial F}{\partial t}=0$，此时固壁上的边界条件为
\begin{equation*}
u\,\frac{\partial F}{\partial x}+v\,\frac{\partial F}{\partial y}+w\,\frac{\partial F}{\partial z}=0
\tag{4.14}
\end{equation*}

(ii) 在固定坐标系中沿$x$正方向匀速移动固壁的边界条件(见图4.1)

此时，在固结于运动物体上的坐标系$(x',y',z')$中壁面方程和时间无关，可写作$F(x',y',z')=0$，因此，在固结于运动物体上的坐标系$(x',y',z')$中，固壁上的边界条件和式(4.14)相同，即
\[
u'\,\frac{\partial F}{\partial x'}+v'\,\frac{\partial F}{\partial y'}+w'\,\frac{\partial F}{\partial z'}=0
\]
式中$u',v',w'$是流体在移动坐标系中的相对速度.

在固定坐标系中匀速移动固壁的边界条件可用以下方法导出.匀速运动坐标系和固定坐标系的变换式为
\[
x'=x-Ut,\qquad y'=y,\qquad z'=z
\]
因而在固定坐标系中界面方程可写作：
\[
F(x-Ut,y,z)=0
\]
此时，$\dfrac{\partial F}{\partial t}=\dfrac{\partial F}{\partial x'}\,\dfrac{\partial x'}{\partial t}=-U\,\dfrac{\partial F}{\partial x'}$，$\dfrac{\partial F}{\partial x}=\dfrac{\partial F}{\partial x'}$，$\dfrac{\partial F}{\partial y}=\dfrac{\partial F}{\partial y'}$，$\dfrac{\partial F}{\partial z}=\dfrac{\partial F}{\partial z'}$，因而在固定坐标系中平移运动固壁上边界条件为
\begin{equation*}
U\,\frac{\partial F}{\partial x}=u\,\frac{\partial F}{\partial x}+v\,\frac{\partial F}{\partial y}+w\,\frac{\partial F}{\partial z}
\tag{4.15}
\end{equation*}
$U$是固壁的移动速度.

(iii) 两种互不掺混液体的可变界面，仍假设界面方程为$F(x,y,z,t)=0$，这时界面两侧流体质点的法向速度应当连续，并等于对应点上界面的法向速度，也就是说界面两侧流体运动都应当满足方程(4.12)，因而有
\begin{equation*}
-\frac{\partial F}{\partial t}=u_{+}\,\frac{\partial F}{\partial x}+v_{+}\,\frac{\partial F}{\partial y}+w_{+}\,\frac{\partial F}{\partial z}=u_{-}\,\frac{\partial F}{\partial x}+v_{-}\,\frac{\partial F}{\partial y}+w_{-}\,\frac{\partial F}{\partial z}
\tag{4.16}
\end{equation*}
下标“$+$”、“$-$”表示界面两侧.

\noindent\textbf{例4.1}\quad 写出流体绕过半径为$R$的固定圆球表面的边界条件.

\noindent\textbf{解}\quad 设圆心在坐标原点，则圆球表面的方程为
\begin{gather*}
F(x,y,z)=x^{2}+y^{2}+z^{2}-R^{2}=0\\[2pt]
\frac{\partial F}{\partial x}=2x\,,\qquad \frac{\partial F}{\partial y}=2y\,,\qquad \frac{\partial F}{\partial z}=2z
\end{gather*}
故圆球上的边界条件为
\[
ux+vy+wz=0
\]
或写成球坐标中的表达式：
\[
u_R=0
\]

\noindent\textbf{例4.2}\quad 写出椭球在流体中沿其长轴方向运动时，椭球面上理想流体运动边界条件，设椭球的三个半轴长度分别为$a,b,c\,(a>b>c)$.

\noindent\textbf{解}\quad 固结在椭球上的物面方程为
\[
\frac{x'^{2}}{a^{2}}+\frac{y'^{2}}{b^{2}}+\frac{z'^{2}}{c^{2}}=1
\]
固结在椭球上的坐标系的边界条件为
\[
\frac{u\,x'}{a^{2}}+\frac{v\,y'}{b^{2}}+\frac{w\,z'}{c^{2}}=0
\]
固结在地面的运动椭球面的方程为
\[
F(x,y,z,t)=\frac{(x-Ut)^{2}}{a^{2}}+\frac{y^{2}}{b^{2}}+\frac{z^{2}}{c^{2}}-1=0
\]
固结在地面的坐标系中，椭球面上流体运动不可穿透条件为
\[
\frac{x-Ut}{a^{2}}\,U=\frac{u(x-Ut)}{a^{2}}+\frac{vy}{b^{2}}+\frac{wz}{c^{2}}
\]

\textcircled{2} 无穷远处条件

当物体在无界静止理想流体中运动时，流体受运动物体驱动，物体对流体所做功总是有限值，因此流体的总动能也应是有限值，对于不可压缩流体距运动物体越远，流体运动速度越小，在无穷远处，流体仍保持静止状态.因而有
\begin{equation*}
|\boldsymbol{x}|\to\infty\,,\quad \boldsymbol{U}=\boldsymbol{0}\,,\quad p=p_\infty\,,\quad \rho=\rho_\infty
\tag{4.17}
\end{equation*}
式中$p_\infty,\rho_\infty$为无穷远处静止流体的热力学状态.

\textcircled{3} 绕流条件

众所周知，运动学边界条件和参考坐标系有关.在静止无界流体中，由匀速运动物体驱动的流场是不定常的，如果把参考坐标系固结在运动物体上，那么物面是静止的，流体由远方流向物体，称这种流动为绕流.由于匀速直线运动的坐标系是惯性坐标系，而在任意惯性坐标系中物体相互作用的动力学关系不变，因此无界静止流体中匀速直线运动物体驱动的流场常用绕流方法处理.设物体以$\boldsymbol{U}_\infty$速度在无界静止的流场中运动，则参考坐标系固结在运动物体上时，无穷远处的来流条件为
\begin{equation*}
|\boldsymbol{x}|\to\infty\,,\quad \boldsymbol{U}=\boldsymbol{U}_\infty\,,\quad p=p_\infty\,,\quad \rho=\rho_\infty
\tag{4.18}
\end{equation*}

(3)两种互不掺混流体界面的应力条件

两种互不掺混流体界面上有流体的表面张力，因此弯曲界面两侧的压强不相等，第1章已经导出弯曲界面两侧压强差为
\begin{equation*}
p_{+}-p_{-}=\gamma\Big(\frac{1}{R_1}+\frac{1}{R_2}\Big)
\tag{4.19}
\end{equation*}
$p_{-}$是凹面上流体压强，$p_{+}$是凸面上流体压强，$\gamma$是表面张力系数.当表面张力可以忽略不计时，界面两侧的压强应当相等.

有了封闭方程组和流动问题的初始及边界条件，用解析方法或是数值方法就可以解出具体的流场.在具体求解流动方程之前，对流动基本性质有足够的了解，会有助于求解具体的流动问题，下面先介绍理想流体在势力场中运动的若干主要性质.

\section*{4.2  理想流体在势力场中运动的主要性质}\addcontentsline{toc}{section}{4.2 理想流体在势力场中运动的主要性质}

\subsection*{1. 凯尔文定理——沿流体线的环量不变定理}

\noindent\textbf{定理4.1}\quad 理想正压流体在势力场中运动时，连续流场内沿封闭流体线的速度环量不随时间变化.

\textbf{先证明一个运动学定理：沿任意封闭流体线速度环量的随体导数等于封闭周线上的加速度环量}，即应有以下等式：
\begin{equation*}
\frac{D}{Dt}\oint_{l(t)}\boldsymbol{U}\cdot\delta\boldsymbol{x}=\oint_{l(t)}\frac{D\boldsymbol{U}}{Dt}\cdot\delta\boldsymbol{x}
\tag{4.20}
\end{equation*}
式(4.20)证明如下：设$l(t)$是$t$时刻的封闭流体线，在$\mathrm{d}t$时刻后它随流体质点移动到$l'(t+\mathrm{d}t)$，考察上述两时刻的速度环量差：
\begin{align*}
\Gamma_{l'}-\Gamma_l&=\oint_{l'}\boldsymbol{U}(\boldsymbol{x}+\mathrm{d}\boldsymbol{x},t+\mathrm{d}t)\cdot\delta\boldsymbol{x}'\\
&\quad-\oint_l\boldsymbol{U}(\boldsymbol{x},t)\cdot\delta\boldsymbol{x}
\end{align*}
为了明确起见，质点位移用$\mathrm{d}\boldsymbol{x}$表示，封闭周线上微元弧长用$\delta\boldsymbol{x}$表示.参考图4.2，流体周线$l'$上微元弧长$\delta\boldsymbol{x}'$和对应流体$l$周线上$\delta\boldsymbol{x}$之间有以下关系：
\begin{align*}
\delta\boldsymbol{x}'&=\delta\boldsymbol{x}+\mathrm{d}\boldsymbol{x}'-\mathrm{d}\boldsymbol{x}\\
&=\delta\boldsymbol{x}+[\boldsymbol{U}'(\boldsymbol{x}+\delta\boldsymbol{x},t)-\boldsymbol{U}(\boldsymbol{x},t)]\,\mathrm{d}t+O(\mathrm{d}t^{2})\\
&=\delta\boldsymbol{x}+\delta\boldsymbol{U}\,\mathrm{d}t+O(\mathrm{d}t^{2})
\end{align*}
式中$\delta\boldsymbol{U}$是$t$时刻在微元线段$\delta\boldsymbol{x}$上两点速度向量差.将此微元弧长的关系式代入前式，可求得$\mathrm{d}t$时间内封闭流体周线上的环量差为
\[
\Gamma_{l'}-\Gamma_l=\oint_{l'}[\boldsymbol{U}(\boldsymbol{x}+\mathrm{d}\boldsymbol{x},t+\mathrm{d}t)-\boldsymbol{U}(\boldsymbol{x},t)]\cdot\delta\boldsymbol{x}+\oint_l\boldsymbol{U}\cdot\delta\boldsymbol{U}\,\mathrm{d}t+O(\mathrm{d}t^{2})
\]
流体线上速度环量的质点导数
\begin{align*}
\frac{\mathrm{d}\Gamma}{\mathrm{d}t}&=\lim_{\mathrm{d}t\to 0}\bigg[\frac{1}{\mathrm{d}t}\Big(\oint_{l'}\boldsymbol{U}'\cdot\delta\boldsymbol{x}'-\oint_l\boldsymbol{U}\cdot\delta\boldsymbol{x}\Big)\bigg]\\
&=\lim_{\mathrm{d}t\to 0}\bigg[\oint\frac{\boldsymbol{U}(\boldsymbol{x}+\mathrm{d}\boldsymbol{x},t+\mathrm{d}t)-\boldsymbol{U}}{\mathrm{d}t}\cdot\delta\boldsymbol{x}+\oint\boldsymbol{U}\cdot\delta\boldsymbol{U}+O(\mathrm{d}t^{2})\bigg]
\end{align*}
因为$\boldsymbol{U}(\boldsymbol{x}+\mathrm{d}\boldsymbol{x},t+\mathrm{d}t)-\boldsymbol{U}(\mathrm{d}\boldsymbol{x},\mathrm{d}t)$是质点速度增量，所以
\[
\lim_{\mathrm{d}t\to 0}\frac{\boldsymbol{U}(\boldsymbol{x}+\mathrm{d}\boldsymbol{x},t+\mathrm{d}t)-\boldsymbol{U}}{\mathrm{d}t}=\boldsymbol{a}
\]
是质点的加速度.另一方面
\[
\oint\boldsymbol{U}\cdot\delta\boldsymbol{U}=\oint\delta(u^{2}+v^{2}+w^{2})
\]
是封闭周线上全微分的积分，故此积分等于零，于是
\[
\frac{\mathrm{d}\Gamma}{\mathrm{d}t}=\oint_l\boldsymbol{a}\cdot\delta\boldsymbol{x}=\oint_l\frac{D\boldsymbol{U}}{Dt}\cdot\delta\boldsymbol{x}
\]
证毕.

下面进一步证明凯尔文定理.理想流体的运动方程为
\[
\frac{D\boldsymbol{U}}{Dt}=-\frac{1}{\rho}\,\nabla p+\boldsymbol{f}
\]
当外力场有势时，质量力是力势$\Pi$的梯度，即$\boldsymbol{f}=-\nabla\Pi$.如流体是正压的，即压强是密度的一元函数，$p=p(\rho)$，则可定义压力函数$\mathscr{P}(\rho)$(见第3章)如下
\[
\mathscr{P}=\int\frac{\mathrm{d}p}{\rho}
\]
或
\begin{gather*}
\mathrm{d}\mathscr{P}=\frac{\mathrm{d}p}{\rho}\\[2pt]
\nabla\mathscr{P}=\frac{\nabla p}{\rho}
\end{gather*}
将质量力势和压力函数代入运动方程，得
\begin{equation*}
\frac{D\boldsymbol{U}}{Dt}=-\nabla\mathscr{P}-\nabla\Pi
\tag{4.21}
\end{equation*}
将上式在封闭流体线上积分，有
\[
\oint\frac{D\boldsymbol{U}}{Dt}\cdot\delta\boldsymbol{x}=-\oint\nabla\mathscr{P}\cdot\delta\boldsymbol{x}-\oint\nabla\Pi\cdot\delta\boldsymbol{x}=0
\]
等式右边两项都是全微分的回路积分，因而等于零.按前面证明的预备定理(式(4.20))，上式左边加速度环量等于速度环量的随体导数，于是有
\begin{equation*}
\frac{\mathrm{d}\Gamma}{\mathrm{d}t}=\oint_l\frac{D\boldsymbol{U}}{Dt}\cdot\delta\boldsymbol{x}=0
\tag{4.22}
\end{equation*}
证毕.

利用凯尔文定理可以推断理想正压流体在势力场中运动时涡量的重要性质.

\subsection*{2. 拉格朗日定理——涡量不生不灭定理}

\noindent\textbf{定理4.2}\quad 理想正压流体在势力场中运动时，如某一时刻连续流场无旋，则流场始终无旋.

\textbf{证明}\quad 设$t=t_0$时刻流场无旋，则流场中任意点都有$\boldsymbol{\omega}=\nabla\times\boldsymbol{U}=\boldsymbol{0}$，因而在流场中通过任意曲面上的涡通量也等于零，即
\[
\int_A\boldsymbol{\omega}\cdot\boldsymbol{n}\,\mathrm{d}A=0
\]
在流体运动学中我们已经证明：任意封闭周线上的速度环量等于通过张在该周线上曲面的涡通量（斯托克斯公式），因此初始时刻任意封闭周线上的环量也等于零，即
\[
\Gamma=\oint_l\boldsymbol{U}\cdot\delta\boldsymbol{x}=\int_A\boldsymbol{\omega}\cdot\boldsymbol{n}\,\mathrm{d}A=0
\]
取任意初始封闭曲线为流体线，当跟随它运动时，根据凯尔文定理，在这些封闭曲线上的环量始终为零.换言之，任意时刻，任意封闭周线上
\[
\Gamma=\oint_{l(t)}\boldsymbol{U}\cdot\delta\boldsymbol{x}=0
\]
再利用斯托克斯公式，可以证明任意时刻通过张在任意封闭曲线$l(t)$上曲面的涡通量等于零，即$\int_{A(t)}\boldsymbol{\omega}\cdot\boldsymbol{n}\,\mathrm{d}A=0$.因为积分式中$A(t)$是流场中的任意曲面（任意位置和任意方向），因而处处有
\[
\boldsymbol{\omega}=\boldsymbol{0}\,,\qquad -\infty<t<\infty
\]
证毕.

拉格朗日定理说明理想正压流体在势力场中运动时涡量是不生不灭的.如果在某一时刻流场无旋则永远无旋，反之如果流场在某一时刻有旋，则永远有旋，其涡量不可能消失.

拉格朗日定理是判断理想正压流体运动是否无旋的理论依据.例如，船舶在静水中开始起航，如果水的粘性可以忽略而且又是不可压缩的，外力只有重力场，则在船舶驱动下，水的运动是无旋的.因为初始时刻水是静止的，处处环量等于零，当然是无旋的，利用拉格朗日定理就可证明水的运动永远无旋.有许多实际流动可近似为无旋流动，如水波流场，飞机以低速直线飞行时的空气运动流场等.

拉格朗日定理也启示我们，流场中涡的产生起因于：①流体的粘性（非理想流体），例如，在静止流体中由物体驱动的流场，在物面有一层很薄的旋涡层；②非正压流场，例如，大气或海洋中的密度分层（非正压）可以形成旋涡；③非有势力场，例如，地球上的气流由于科氏力（非有势力）的作用可以生成旋涡，在极端情况下，可以形成很强烈的旋风；④流场的间断（非连续流场）也可能生成旋涡，例如，在高速气流中的曲面激波后，就会产生有旋流动.

\subsection*{3. 亥姆霍兹定理——涡线及涡管保持定理}

\textbf{(1)涡线保持定理（亥姆霍兹第一定理）}

一般情况涡线可视作两涡面的交线.因此，只要能证明涡面有保持性，就证明了涡线的保持性.所谓涡面保持性，就是说：组成涡面的流体质点永远组成涡面，或者说涡面是流体面.

\noindent\textbf{定理4.3}\quad 理想正压流体在势力场中运动时，组成涡面的流体质点永远组成涡面.

\textbf{证明}\quad 设某一时刻，流体中有一涡面$\Sigma$，在涡面上有$\omega_n=\boldsymbol{\omega}\cdot\boldsymbol{n}=0$.在此涡面上，任取一封闭周线$l$，则由斯托克斯定理可以证明，绕此封闭曲线的环量等于零.原封闭曲线$l$随曲面运动到$\Sigma'$时变形为$l'$，它是流体线.根据凯尔文定理，在任意时刻封闭流体线$l'$上的环量为零，也就是说，曲面$\Sigma'$上任意封闭周线的环量等于零，因此$\Sigma'$面上处处有$\omega_n=\boldsymbol{\omega}\cdot\boldsymbol{n}=0$.于是$\Sigma'$也是涡面.涡面保持定理说明，理想正压流体在势力场运动时，涡面始终是涡面.下面进一步证明涡线保持定理.

\noindent\textbf{定理4.4}\quad 理想正压流体在势力场中运动时，组成涡线的质点永远组成涡线.

\textbf{证明}\quad 设某一时刻有一涡线$l$，它是两个涡面$\Sigma_1$和$\Sigma_2$的交线.涡面保持定理告诉我们，这两个涡面在任意时刻始终是由相同质点组成的涡面$\Sigma_1(t)$和$\Sigma_2(t)$，因此它们的交线始终是由相同质点组成的涡线.孤立涡线（即涡线以外都是无旋的）也可以用类似方法证明：设有初始两个流体面的交线是给定涡线，它始终是流体面的交线，而流体面其他地方涡量总是等于零，因此孤立涡线始终是孤立涡线.

涡线保持定理说明，理想正压流体在势力场中运动时，涡线始终是流体线.

最后证明涡管和涡管强度保持性.

\noindent\textbf{定理4.5}\quad 理想正压流体在势力场中运动时，组成涡管的流体质点始终组成涡管，并且涡管的强度不随时间变化（亥姆霍兹第二定理）.

\textbf{证明}\quad 涡管表面是涡面，因此由涡面保持性定理就证明了涡管的保持性.就是说涡管表面始终是流体面，因此涡管内的流体质点永远在涡管内部.进一步证明涡管强度保持性.在涡管表面上作围绕涡管的任意封闭曲线$l$，沿此封闭曲线的环量等于同一瞬间的涡管通量，也就是该涡管的强度.上述封闭曲线随涡管运动时是流体线并且始终围绕涡管，根据凯尔文定理，沿此封闭线的环量是不变的，因此涡管的涡通量，即它的强度，不随时间变化，于是证明了涡管强度的保持性.

涡管和涡管强度保持性说明，理想正压流体在势力场中运动时，涡管和涡线或涡面一样，它在运动过程中可以变形，但是组成涡管的质点不变，涡管的强度也不变.

凯尔文定理、拉格朗日定理和亥姆霍兹定理全面地描述了理想正压流体在势力场中运动时，涡量演化的基本规律.简要地说，理想正压流体在势力场作用下运动时，无旋运动永远是无旋的；有旋运动永远是有旋的；涡管始终是涡管，并且它的强度不变.

本章后面几节主要讨论无旋流动，流体的有旋运动将在本章最后讨论.

\section*{4.3  兰姆型方程和理想流体运动的几个积分}\addcontentsline{toc}{section}{4.3 兰姆型方程和理想流体运动的几个积分}

\subsection*{1. 兰姆（Lamb）型方程}

欧拉运动方程为
\[
\frac{\partial\boldsymbol{U}}{\partial t}+\boldsymbol{U}\cdot\nabla\boldsymbol{U}=-\frac{1}{\rho}\,\nabla p+\boldsymbol{f}
\]
应用向量导数运算公式
\[
\boldsymbol{U}\cdot\nabla\boldsymbol{U}=\frac{1}{2}\,\nabla(\boldsymbol{U}\cdot\boldsymbol{U})-\boldsymbol{U}\times(\nabla\times\boldsymbol{U})
\]
将上式替换欧拉方程中的对流导数项，理想流体运动方程可写成
\begin{equation*}
\frac{\partial\boldsymbol{U}}{\partial t}+\frac{1}{2}\,\nabla(\boldsymbol{U}\cdot\boldsymbol{U})-\boldsymbol{U}\times\boldsymbol{\omega}=-\frac{1}{\rho}\,\nabla p+\boldsymbol{f}
\tag{4.23}
\end{equation*}
式(4.23)称为兰姆型方程.

\subsection*{2. 理想正压流体在势力场中运动的两个积分式}

\textbf{(1)伯努利积分}：理想正压流体在势力场中作定常流动时沿流线有
\begin{equation*}
\frac{1}{2}\,|\boldsymbol{U}|^{2}+\mathscr{P}+\Pi=C(n)
\tag{4.24}
\end{equation*}
式中$n$代表不同流线.

\textbf{证明}\quad 定常流动中$\dfrac{\partial\boldsymbol{U}}{\partial t}=\boldsymbol{0}$，正压流体有压强函数$\mathscr{P}=\displaystyle\int\frac{\mathrm{d}p}{\rho}$（式(3.42)），于是运动方程(4.23)可写成
\[
\nabla\Big(\frac{1}{2}\,|\boldsymbol{U}|^{2}+\mathscr{P}+\Pi\Big)-\boldsymbol{U}\times\boldsymbol{\omega}=\boldsymbol{0}
\]
将方程在流线的切线方向$\boldsymbol{s}$上投影，得
\[
\boldsymbol{s}\cdot\nabla\Big(\frac{1}{2}\,|\boldsymbol{U}|^{2}+\mathscr{P}+\Pi\Big)-\boldsymbol{s}\cdot(\boldsymbol{U}\times\boldsymbol{\omega})=0
\]
由于$(\boldsymbol{U}\times\boldsymbol{\omega})\perp\boldsymbol{U}$，而$\boldsymbol{s}\,\|\,\boldsymbol{U}$，故最后一项等于零.另外，根据梯度定义：$\boldsymbol{s}\cdot\nabla=\dfrac{\partial}{\partial s}$是沿流线方向的方向导数，于是有
\[
\frac{\partial}{\partial s}\Big(\frac{1}{2}\,|\boldsymbol{U}|^{2}+\mathscr{P}+\Pi\Big)=0
\]
积分后得
\[
\frac{1}{2}\,|\boldsymbol{U}|^{2}+\mathscr{P}+\Pi=C(n)
\]
证毕.

伯努利积分（也称伯努利公式）中$C(n)$是同一流线上的积分常数，但在不同流线上（指标$n$变化）积分常数$C(n)$可以变化，$C(n)$称为伯努利常数.

对于不可压缩流体，$\rho=$常数，故$\mathscr{P}=\displaystyle\int\frac{\mathrm{d}p}{\rho}=\frac{p}{\rho}$，沿流线的伯努利公式可简化为
\begin{equation*}
\frac{1}{2}\,|\boldsymbol{U}|^{2}+\frac{p}{\rho}+\Pi=C(n)
\tag{4.25}
\end{equation*}
伯努利常数$C(n)$可由流线起始点上的参数给出，若流场的起始面（或不是流面的任意曲面）上流动均匀，该面上伯努利常数处处相等，则全场有$C(n)=\mathrm{const.}$

\textbf{(2)}理想正压流体在势力场中作定常流动时，沿涡线有伯努利积分
\begin{equation*}
\frac{1}{2}\,|\boldsymbol{U}|^{2}+\mathscr{P}+\Pi=C(m)
\tag{4.26}
\end{equation*}
将运动方程沿涡线积分便可得式(4.26)，证明方法同流线上的伯努利积分，读者可自己证明.式(4.26)中积分常数$C(m)$在同一涡线上不变.

\subsection*{3. 柯西-拉格朗日积分}

理想正压流体在势力场作无旋流动时，全场有下式成立
\begin{equation*}
\frac{\partial\Phi}{\partial t}+\frac{1}{2}\,|\nabla\Phi|^{2}+\mathscr{P}+\Pi=C(t)
\tag{4.27a}
\end{equation*}
\textbf{证明}\quad 当流动无旋时，有速度势$\boldsymbol{U}=\nabla\Phi$，及$\boldsymbol{\omega}=\nabla\times\boldsymbol{U}=\boldsymbol{0}$，将速度势代入兰姆型方程得
\[
\nabla\frac{\partial\Phi}{\partial t}+\nabla\frac{1}{2}\,|\nabla\Phi|^{2}+\nabla\mathscr{P}+\nabla\Pi=\boldsymbol{0}
\]
上式可写作
\[
\nabla\Big(\frac{\partial\Phi}{\partial t}+\frac{1}{2}\,|\nabla\Phi|^{2}+\mathscr{P}+\Pi\Big)=\boldsymbol{0}
\]
全场积分得
\[
\frac{\partial\Phi}{\partial t}+\frac{1}{2}\,|\nabla\Phi|^{2}+\mathscr{P}+\Pi=C(t)
\]
积分常数$C(t)$仅和时间有关，与空间坐标无关，也就是说，同一时刻在所有流线上积分常数是相同的.对于定常无旋流动，则积分常数$C(t)$和时间无关.换句话说，理想正压流体在势力场中作定常无旋流动时，伯努利积分常数在全流场相同，即有
\begin{equation*}
\frac{1}{2}\,|\nabla\Phi|^{2}+\mathscr{P}+\Pi=C
\tag{4.27b}
\end{equation*}

伯努利积分是理想流体运动中常用的公式.第3章中在流动的一维近似假定下，通过积分型动量方程曾导出理想正压流体在势力场中定常运动时沿流管的伯努利积分，它是一般伯努利公式的特例和近似.

应当注意不同陈述方式的伯努利公式之间的差别.它们共同的条件是：理想不可压缩流体在势力场中运动；它们的不同之处是：全流场伯努利常数相同的公式只适用于无旋流动；沿流线或流管伯努利常数相等的公式既适用于无旋流动，也适用于有旋流动.

拉格朗日定理证明了理想正压流体在势力场的无旋流动始终保持无旋流动，许多理想流体运动可以近似为无旋流动.下面研究无旋流动的性质.

\section*{4.4  理想不可压缩无旋流动问题的数学提法及主要性质}\addcontentsline{toc}{section}{4.4 理想不可压缩无旋流动问题的数学提法及主要性质}

不可压缩流体空间无旋流动问题的数学提法以位势场论为基础，下面介绍这类流动问题的理论和分析计算方法.

\subsection*{1. 不可压缩无旋流动的基本方程}

不可压缩流场中速度场的散度等于零
\begin{equation*}
\nabla\cdot\boldsymbol{U}=0
\tag{4.28}
\end{equation*}
而无旋流场中速度场的旋度应等于零，即
\begin{equation*}
\nabla\times\boldsymbol{U}=\boldsymbol{0}
\tag{4.29}
\end{equation*}
因而有速度势：
\begin{equation*}
\boldsymbol{U}=\nabla\Phi
\tag{4.30}
\end{equation*}
将式(4.30)代入式(4.28)得
\begin{equation*}
\Delta\Phi=0
\tag{4.31}
\end{equation*}
即不可压缩无旋流动的速度势满足拉普拉斯方程.

\subsection*{2. 边界条件}

\textbf{(1)固壁运动学条件}

设已知壁面方程为$F(x,y,z,t)=0$，则由4.1节导出的固壁不可穿透条件为
\[
\frac{\partial F}{\partial t}+u\,\frac{\partial F}{\partial x}+v\,\frac{\partial F}{\partial y}+w\,\frac{\partial F}{\partial z}=0
\]
或
\begin{equation*}
\frac{\partial F}{\partial t}+\frac{\partial\Phi}{\partial x}\,\frac{\partial F}{\partial x}+\frac{\partial\Phi}{\partial y}\,\frac{\partial F}{\partial y}+\frac{\partial\Phi}{\partial z}\,\frac{\partial F}{\partial z}=0
\tag{4.32}
\end{equation*}

\textbf{(2)无穷远处渐近条件}

对于绕流问题
\begin{equation*}
|\boldsymbol{x}|\to\infty:\ \boldsymbol{U}=\nabla\Phi=\boldsymbol{U}_\infty\,,\qquad p=p_\infty\,,\qquad \rho=\rho_\infty
\tag{4.33a}
\end{equation*}
物体在静止流体中运动时，无穷远处流体的渐近条件：
\begin{equation*}
|\boldsymbol{x}|\to\infty:\ \boldsymbol{U}=\nabla\Phi=\boldsymbol{0}\,,\qquad p=p_\infty
\tag{4.33b}
\end{equation*}

下面将说明在已知界面运动规律的无界单连域流场中，给定运动学边界条件式(4.32)和式(4.33)，速度势方程(4.31)是定解的.也就是说，单连域中不可压缩无旋流动的速度势完全由运动学边界条件确定，或者说，速度场完全由运动学条件确定.

\subsection*{3. 压强分布}

速度场确定以后，由柯西-拉格朗日积分$\dfrac{\partial\Phi}{\partial t}+\dfrac{1}{2}\,|\nabla\Phi|^{2}+\dfrac{p}{\rho}+\Pi=C(t)$确定压强分布，其中积分函数$C(t)$由边界条件确定.例如可以用无穷远处一点的参数确定$C(t)$：
\begin{equation*}
C(t)=\Big(\frac{\partial\Phi}{\partial t}\Big)_{\infty}+\frac{1}{2}\,\Big|\nabla\Phi\Big|_{\infty}^{2}+\frac{p_\infty}{\rho}+\Pi_\infty
\tag{4.34}
\end{equation*}
以上分析说明了求解给定边界条件的不可压缩无旋流动的步骤是：

第一步：由运动学条件求速度势，即速度场；

第二步：由柯西-拉格朗日积分求压强分布.

\subsection*{4. 应用举例}

\noindent\textbf{例4.3}\quad 圆球在无界静止的理想不可压缩流体中匀速膨胀$R_a=R_a(t)$.设质量力为零，求圆球外的速度场和压强分布（$t=0$时，球面速度为零）.

\noindent\textbf{解}\quad 因流动由静止启动，根据拉格朗日定理判断流动是无旋的，基本方程是
\[
\Delta\Phi=0\,,\qquad \frac{\partial\Phi}{\partial t}+\frac{1}{2}\,|\nabla\Phi|^{2}+\frac{p}{\rho}=C(t)
\]
边界条件为：在球面$x^{2}+y^{2}+z^{2}-R_a^{2}=0$上不可穿透条件：
\[
(\nabla\Phi)\cdot\boldsymbol{n}=\boldsymbol{U}_b\cdot\boldsymbol{n}=\dot{R}_a(t)
\]
在球坐标中，它可写作：
\[
\Big(\frac{\partial\Phi}{\partial R}\Big)_{R=R_a}=\dot{R}_a(t)
\]
无穷远处静止条件：
\[
|\boldsymbol{x}|\to\infty:\ \nabla\Phi=\boldsymbol{0}\,,\qquad p=p_0
\]
本题给定的边界条件是球对称的，根据拉普拉斯方程的性质，其解也应是球对称的，即：$\Phi=\Phi(R,t)$.$R$为球坐标的向径，球对称条件下拉普拉斯方程可表示为
\[
\frac{\partial}{\partial R}\Big(R^{2}\,\frac{\partial\Phi}{\partial R}\Big)=0
\]
它的一般解为
\[
\Phi=\frac{C_1(t)}{R}+C_2(t)
\]
将边界条件$\Big(\dfrac{\partial\Phi}{\partial R}\Big)_{R=R_a}=\dot{R}_a(t)$代入可得
\[
C_1=-R_a^{2}(t)\,\dot{R}_a(t)
\]
$C_2(t)$并不影响所解流场，可以取为零，从而
\[
\Phi=\frac{-R_a^{2}(t)\,\dot{R}_a(t)}{R}
\]
由速度势梯度计算出速度场等于：
\[
\boldsymbol{U}=\nabla\Phi=\frac{\partial}{\partial R}\Big(\frac{-R_a^{2}(t)\,\dot{R}_a(t)}{R}\Big)\boldsymbol{e}_R=\frac{R_a^{2}(t)\,\dot{R}_a(t)}{R^{2}}\,\boldsymbol{e}_R
\]
它满足无穷远处的静止条件：
\[
|\boldsymbol{U}|_{R\to\infty}=|\nabla\Phi|_{R\to\infty}=\Big(\frac{R_a^{2}(t)\,\dot{R}_a(t)}{R^{2}}\Big)_{R\to\infty}=0
\]
最后，应用柯西-拉格朗日积分式求压强分布，先由无穷远处条件求积分常数$C(t)$，因为
\[
\Big(\frac{\partial\Phi}{\partial t}\Big)_{R\to\infty}=-\Big(\frac{2\dot{R}_a^{2}(t)\,R_a(t)}{R}+\frac{\ddot{R}_a(t)\,R_a^{2}(t)}{R}\Big)_{R\to\infty}=0
\]
\[
|\nabla\Phi|_{R\to\infty}=\Big(\frac{R_a^{2}(t)\,\dot{R}_a(t)}{R^{2}}\Big)_{R\to\infty}=0
\]
代入式(4.34)，可得：$C(t)=\dfrac{p_0}{\rho}$.将积分常数代入柯西-拉格朗日积分后，得流场的压强分布等于：
\begin{align*}
p&=p_0-\rho\Big(\frac{\partial\Phi}{\partial t}+\frac{1}{2}\,|\nabla\Phi|^{2}\Big)\\
&=p_0+\rho\bigg[\frac{2\dot{R}_a^{2}(t)\,R_a(t)}{R}+\frac{\ddot{R}_a(t)\,R_a^{2}(t)}{R}\bigg]-\frac{\rho}{2}\,\frac{\dot{R}_a^{2}(t)\,R_a^{4}(t)}{R^{4}}
\end{align*}
当球等速膨胀时$\ddot{R}_a=0$，
\[
p=p_0+\frac{2\,\rho R_a\dot{R}_a^{2}}{R}-\frac{\rho}{2}\,\frac{R_a^{4}\dot{R}_a^{2}}{R^{4}}
\]
球面上的压强为
\[
p_{R=R_a(t)}=p_0+\frac{3}{2}\,\rho\dot{R}_a^{2}
\]

\subsection*{5. 速度势$\Phi$的物理意义}

\noindent\textbf{定义4.1}\quad 压强对时间的积分称为压强冲量，并用$I$表示，即$I=\displaystyle\int_0^{\tau}p\,\mathrm{d}t$.

可以证明：静止不可压缩理想流体在瞬时脉冲压强作用下产生的流动是无旋的，它的速度势等于负压强冲量除以密度.

设想不可压缩理想流体的静止流场各点在极短时间内（$\delta t\ll 1$）同时受到连续分布的压强冲量作用而启动.可以证明此时产生的流场是无旋的，且速度势$\Phi$等于负压强冲量除以密度，即$\Phi=-\dfrac{I}{\rho}$.在$\delta t$时间内对欧拉方程进行积分：
\[
\int_0^{\delta t}\Big(\frac{\partial\boldsymbol{U}}{\partial t}+\boldsymbol{U}\cdot\nabla\boldsymbol{U}\Big)\mathrm{d}t=\int_0^{\delta t}\Big(-\frac{\nabla p}{\rho}+\boldsymbol{f}\Big)\mathrm{d}t
\]
因初始流场静止：$\boldsymbol{U}(\boldsymbol{x},0)=\boldsymbol{0}$，左端第一项积分$\displaystyle\int_0^{\delta t}\frac{\partial\boldsymbol{U}}{\partial t}\,\mathrm{d}t=\boldsymbol{U}(\boldsymbol{x},\delta t)-\boldsymbol{U}(\boldsymbol{x},0)=\boldsymbol{U}(\boldsymbol{x},\delta t)$；

左端第二项积分$\displaystyle\int_0^{\delta t}\boldsymbol{U}\cdot\nabla\boldsymbol{U}\,\mathrm{d}t=O(\delta t)$；

右端第一项积分$-\displaystyle\int_0^{\delta t}\frac{\nabla p}{\rho}\,\mathrm{d}t=-\frac{1}{\rho}\,\nabla\int_0^{\delta t}p\,\mathrm{d}t=-\frac{\nabla I}{\rho}$；

右端第二项积分$\displaystyle\int_0^{\delta t}\boldsymbol{f}\,\mathrm{d}t=O(\delta t)$.

将以上积分代入欧拉方程的积分式，并令$\delta t\to 0$，得
\[
\boldsymbol{U}=-\frac{\nabla I}{\rho}=-\nabla\,\frac{I}{\rho}
\]
上式表明该速度场是某一函数的梯度，因而是无旋的，且有$\boldsymbol{U}=\nabla\Phi$，其中
\[
\Phi=-\frac{I}{\rho}+C
\]
常数$C$不影响速度场，可以取为零.于是证明了速度势等于负瞬时压强冲量除以密度.

速度势的物理意义表明不可压缩流体的无旋流动可由瞬时压强冲量产生.事实上，即使流体有粘性，例如欧拉方程右端加上单位质量的粘性力，在极短时间内粘性力的积分也趋于零.就是说由瞬时压强冲量产生的不可压缩流动，在初始瞬间是无旋的.如果流体的粘性可以忽略，则初始的无旋流动保持无旋；如果流体有粘性则初始无旋流动因粘性而演化为有旋流动.由于瞬时压强冲量正比于速度势，因此它也满足拉普拉斯方程，这时给定边界上瞬时压强冲量或它的法向导数就能确定域内的压强冲量，换句话说瞬时的速度势边值就确定该瞬时域内压强冲量及该瞬时的流场.例如，刚性物体在无界静止的不可压缩流体中突然启动，根据上面的分析，启动瞬间的流场是无旋的，这时物面上的法向速度等于当地的物体速度的法向分量，也就是已知物面上的$\partial\Phi/\partial n=-\partial(I/\rho)/\partial n$.由于不可压缩流体无旋运动的速度势满足拉普拉斯方程，当给定边界上的速度势法向导数后，流场的速度势，也就是瞬时压强脉冲除以密度就确定了（不计任意常数），于是，瞬时速度场就唯一地确定.\textbf{对于不可压缩流体的无旋流动，瞬时的速度场由该瞬时的边界法向速度唯一地确定，而和流场过去的历史无关，这是不可压缩流体无旋流动的重要性质.}

\subsection*{6. 不可压缩无旋流场的主要特性}

可以证明不可压缩无旋流场有以下特性.

\textbf{(1)速度势函数不能在域内有极大或极小值}

\textbf{证明}\quad 不可压缩流体流动有以下的连续方程：
\[
\nabla\cdot\boldsymbol{U}=0
\]
因而在域内任意封闭曲面上的体积流量等于零，即：$\displaystyle\oint_\Sigma\boldsymbol{n}\cdot\boldsymbol{U}\,\mathrm{d}A=\int_V\nabla\cdot\boldsymbol{U}\,\mathrm{d}V=0$，无旋流场中：$\boldsymbol{U}=\nabla\Phi$，故有：$\displaystyle\oint_\Sigma\frac{\partial\Phi}{\partial n}\,\mathrm{d}A=0$.可以用反证法来证明$\Phi$的极值性质.设在域内某点$\Phi$为极大值，则在该极值点邻域中的一小球上处处有$\dfrac{\partial\Phi}{\partial n}<0$，因而$\displaystyle\oint_\Sigma\frac{\partial\Phi}{\partial n}\,\mathrm{d}A<0$，这和速度势基本性质$\displaystyle\oint_\Sigma\frac{\partial\Phi}{\partial n}\,\mathrm{d}A=0$相违背.同理可以证明$\Phi$不能在域内有极小值，于是定理得证.

\textbf{(2)不可压缩无旋流场中速度的模不能在流场中达到极大值}

\textbf{证明}\quad 仍用反证法.设流场中某点$A$的速度模达到极大值$U_A$，且$U_A$的方向为坐标轴$x$正方向，即$U_A=\Big(\dfrac{\partial\Phi}{\partial x}\Big)_A$.速度势$\Phi$满足拉普拉斯方程$\nabla\Phi=0$，因此$\dfrac{\partial\Phi}{\partial x}$也满足拉普拉斯方程.根据上面已证明的性质，$\dfrac{\partial\Phi}{\partial x}$不可能在域内有极值，因此在$A$点附近必能找到一点$M$，该处$\Big(\dfrac{\partial\Phi}{\partial x}\Big)_M>\Big(\dfrac{\partial\Phi}{\partial x}\Big)_A$.该处速度
\[
U_M=\sqrt{\Big(\frac{\partial\Phi}{\partial x}\Big)_M^{2}+\Big(\frac{\partial\Phi}{\partial y}\Big)_M^{2}+\Big(\frac{\partial\Phi}{\partial z}\Big)_M^{2}}\geqslant\Big(\frac{\partial\Phi}{\partial x}\Big)_M>U_A
\]
这和假定$U_A$是极大值相违背，因而$U_A$不能在流场中取极大值.

速度值的极值定理说明，\textbf{不可压缩无旋流场中速度最大值必在物面或其他流动界面上}.

\textbf{(3)重力场中不可压缩无旋流压强不能在域内达到极小值}

\textbf{证明}\quad 在$\Delta\Phi=0$的势流场中，可以导出域内有$\displaystyle\oint_\Sigma\frac{\partial p}{\partial n}\,\mathrm{d}A<0$.由不可压缩无旋流场中的柯西-拉格朗日积分：$\dfrac{\partial\Phi}{\partial t}+\dfrac{|\boldsymbol{U}|^{2}}{2}+\dfrac{p}{\rho}+gz=0$，将拉普拉斯算符作用于等式两边，因为$\Delta\Phi=0$和$\Delta z=0$，得
\begin{equation*}
\Delta p=-\frac{1}{2}\,\rho\Delta\,|\boldsymbol{U}|^{2}
\tag{4.35}
\end{equation*}
而
\begin{align*}
\Delta\,|\boldsymbol{U}|^{2}&=\Delta u^{2}+\Delta v^{2}+\Delta w^{2}\\
&=2u\Delta u+2v\Delta v+2w\Delta w+2\big(\,|\nabla u\,|^{2}+|\nabla v\,|^{2}+|\nabla w\,|^{2}\big)
\end{align*}
其中
\[
\Delta u=\Delta\Big(\frac{\partial\Phi}{\partial x}\Big)=\frac{\partial}{\partial x}(\Delta\Phi)=0
\]
同理$\Delta v=\Delta w=0$，因而
\[
\Delta p=-\frac{\rho}{2}\big(\,|\nabla u\,|^{2}+|\nabla v\,|^{2}+|\nabla w\,|^{2}\big)\leqslant 0
\]
利用高斯公式，有
\begin{equation*}
\oint_\Sigma\frac{\partial p}{\partial n}\,\mathrm{d}S=\int_V\Delta p\,\mathrm{d}V=-\frac{\rho}{2}\int_V\big(\,|\nabla u\,|^{2}+|\nabla v\,|^{2}+|\nabla w\,|^{2}\big)\mathrm{d}V<0
\tag{4.36}
\end{equation*}
如果域内有一点压强为极小值，则在围绕该点的邻域的一个小球面上处处有$\dfrac{\partial p}{\partial n}>0$，因而该面上$\displaystyle\oint_\Sigma\frac{\partial p}{\partial n}\,\mathrm{d}S>0$，这和式(4.36)相违背.因此域内压强不能达到极小值.换言之，\textbf{不可压缩无旋流动中压强的极小值只能在物面上}.

不可压缩无旋流动中的压强极值定理，不仅有理论意义，也有实用意义.例如在正常情况下，水中绕流物体上的压强应大于当地的汽化压强，否则在物面上就会发生汽化空泡，既影响绕流物面上的作用力，又可能发生使物面毁坏的气蚀现象.根据压强极值定理，汽化首先发生在物面上，不可能发生在流场内部.为了避免发生气蚀现象，应当注意水中绕流物体的外形设计，使得物面上的最小压强不要太低.通常物面上的凸起或尖角处的压强很低，所以应当使绕流物面十分光滑，以减小发生汽蚀的可能性.

\subsection*{7. 不可压缩无旋流动中速度势的数学性质}

前面已经讨论了理想不可压缩流体在有势力场中的无旋运动的求解方法，分析计算这类问题的关键是求解速度势的拉普拉斯方程.为此先介绍拉普拉斯方程定解的理论，这部分内容在数理方程课程中有详细论述，本书只援引结论而不加证明.后续各节将应用这些数学方法求解具体流动问题.首先介绍求解域的几何性质.

\textbf{(1)单连通域和多连通域}

\noindent\textbf{定义4.2}\quad 任意封闭周线可以在域内收缩到一点的空间域为单连通域，否则为多连通域.

图4.3中列举的三个几何图形中，圆球外是单连通域，球内也是单连通域，因为圆球内外的任意周线都是可缩的；无限长圆柱体内域是单连通域，外域是双连通域，因为围绕圆柱的任意周线，不可能在圆柱外收缩到一点.圆环的内域和外域都是双连通域，因为圆环外绕圆环的周线是不可缩的，圆环内的周线（虚线）也是不可缩的.

多连通域的性质可以用如下的分隔面来分类.

\noindent\textbf{定义4.3}\quad 完全在域内并和域的封闭边界相交的曲面称为隔面（参见图4.4）.

\noindent\textbf{定义4.4}\quad 需加$n-1$个隔面而不破坏连通性并使之成为单连通域的空间域称为$n$连通域.

该定义表明，加一个隔面而不破坏连通性的空间域是双连通域；加两个隔面而不破坏连通性的空间域是三连通域.例如：圆环内域（图4.4(a)）用一回截面分隔后仍保持原有连通性，但不能再加任何分隔面而不破坏环内的连通性，因此圆环内域是双连通域.用一无限大半平面作分割面和圆柱体相交于母线（图4.4(b)），这时分割后的无穷长圆柱外域也是单连通域.

\textbf{(2)双连通域中速度势的多值性}

无旋速度场的速度势可以用以下线积分确定：
\[
\Phi=\int_{x_0}^{x}\boldsymbol{U}\cdot\mathrm{d}\boldsymbol{x}
\]
速度势可能存在多值性，它由速度环量$\displaystyle\oint_l\boldsymbol{U}\cdot\mathrm{d}\boldsymbol{x}$的积分来确定.因为速度势的线积分公式可以通过无限条路径来实现，所以速度势的一般公式可写成
\begin{equation*}
\Phi(\boldsymbol{x})=\Phi_P(\boldsymbol{x})+\oint_l\boldsymbol{U}\cdot\mathrm{d}\boldsymbol{x}
\tag{4.37}
\end{equation*}
式中$\Phi_P(\boldsymbol{x})$是速度势积分式中某一单向积分路径的积分值，$l$是通过$\boldsymbol{x}$和$\boldsymbol{x}_0$的任意封闭曲线.如果式(4.37)中回路积分不等于零，则同一点的速度势可以有无穷多值.由上式可见：若流场中任意封闭周线上速度环量都等于零，则速度势是单值的，于是有以下结论.

\textcircled{1} \textbf{在单连通域中无旋流动的速度势是单值的.}

因为在单连通的无旋流场中，任何封闭周线都是可缩的，并且任何封闭周线上所张的曲面$\Sigma$都在域内，而在无旋流场内，$\nabla\times\boldsymbol{U}=\boldsymbol{0}$，应用斯托克斯公式，得
\[
\oint_l\boldsymbol{U}\cdot\mathrm{d}\boldsymbol{x}=\oint_\Sigma(\nabla\times\boldsymbol{U})\cdot\boldsymbol{n}\,\mathrm{d}A=0
\]
即任意封闭周线上$\displaystyle\oint_l\mathrm{d}\Phi=\oint_l\boldsymbol{U}\cdot\mathrm{d}\boldsymbol{x}=0$，因此速度势是单值的.

\textcircled{2} \textbf{双连通域的无旋流场中，任意不可缩周线上的速度环量（只绕封闭周线一周）相等.}

现在考察双连通域中的无旋运动，这时在可缩周线上，仍有
\[
\oint_l\boldsymbol{U}\cdot\mathrm{d}\boldsymbol{x}=\oint_\Sigma(\nabla\times\boldsymbol{U})\cdot\boldsymbol{n}\,\mathrm{d}A=0
\]
但在不可缩周线$L$上，积分$\displaystyle\oint_L\boldsymbol{U}\cdot\mathrm{d}\boldsymbol{x}$的值不能确定，因为$l$所包围的空间并不都是无旋流场，总有一部分不属于无旋流场，因此在双连通域中无旋流动的速度势有可能是多值的.但双连通域中的环量积分式有下列性质：\textbf{任意不可缩周线上的速度环量相等}.

\textbf{证明}\quad 在图4.5所示的双连通域中任取两个不可缩周线$L_\Gamma$和$L_{\Gamma'}$（$L_\Gamma$和$L_{\Gamma'}$表示逆时针周线，$L_{-\Gamma}$和$L_{-\Gamma'}$表示顺时针周线），为要证明
\[
\oint_{L_\Gamma}\boldsymbol{U}\cdot\mathrm{d}\boldsymbol{x}=\oint_{L_{\Gamma'}}\boldsymbol{U}\cdot\mathrm{d}\boldsymbol{x}
\]
在$L_\Gamma$和$L_{\Gamma'}$之间作一分隔线（是分隔面在平面上的交线）.为了证明方便，把隔线作成两无限靠近的线段$AB$、$CD$.现在由$L_\Gamma+AB+L_{-\Gamma'}+CD$组成的周线在域内是可缩的，因此
\[
\oint_{L_\Gamma+AB+L_{-\Gamma'}+CD}\boldsymbol{U}\cdot\mathrm{d}\boldsymbol{x}=0
\]
由于$AB$和$CD$方向相反，故：$\displaystyle\int_{AB+CD}\boldsymbol{U}\cdot\mathrm{d}\boldsymbol{x}=0$，从而$\displaystyle\int_{L_\Gamma+L_{-\Gamma'}}\boldsymbol{U}\cdot\mathrm{d}\boldsymbol{x}=0$，或
\[
\oint_{L_\Gamma}\boldsymbol{U}\cdot\mathrm{d}\boldsymbol{x}+\oint_{L_{-\Gamma'}}\boldsymbol{U}\cdot\mathrm{d}\boldsymbol{x}=0
\]
因为$L_{-\Gamma'}$是顺时针方向的周线，故有$\displaystyle\oint_{L_{-\Gamma'}}\boldsymbol{U}\cdot\mathrm{d}\boldsymbol{x}=-\oint_{L_{\Gamma}}\boldsymbol{U}\cdot\mathrm{d}\boldsymbol{x}$，最后得
\[
\oint_{L_\Gamma}\boldsymbol{U}\cdot\mathrm{d}\boldsymbol{x}=\oint_{L_{\Gamma'}}\boldsymbol{U}\cdot\mathrm{d}\boldsymbol{x}
\]
于是证明了双连通域中不可缩周线上的环量是常数，并令$\displaystyle\oint_L\boldsymbol{U}\cdot\mathrm{d}\boldsymbol{x}=\Gamma$.已知双连通域中不可缩周线上的环量后，速度势的多值性可由不可缩周线上的环量来确定.当积分路径单向由$\boldsymbol{x}_0$到$\boldsymbol{x}$时，$\displaystyle\int_{x_0}^{x}\boldsymbol{U}\cdot\mathrm{d}\boldsymbol{x}=\Phi_P$；如果积分路径由$\boldsymbol{x}_0$出发绕不可缩周线$n$圈后到达$\boldsymbol{x}$，则积分值等于
\begin{equation*}
\Phi=\Phi_P+n\Gamma
\tag{4.38}
\end{equation*}
式中$n\geqslant 1$是正整数，式(4.38)说明了双连通域中的速度势的多值性.如果双连通域中的无旋流动绕不可缩周线的环量等于零，那么速度势仍是单值的.

以上的数学理论告诉我们：单连通域中的无旋流动是单值有势流动；多连通域中的无旋流动可能是有环量的有势流动.

双连通域中存在有环量的无旋流动，可从运动学上作如下解释，设流场中有一无限长的涡管，涡管以外处处无旋，围绕涡管周线的环量处处等于涡管强度.无限长涡管以外的空间是双连通域，所以这是一个典型的双连通域中有环量的无旋流动.又如：绕无限长机翼的翼型截面的平面流动也是典型的双连通域的流动现象，这种流动现象可近似为有环量的无旋流动.绕机翼环量可以设想为无限长机翼内部有一根或若干根无限长的涡管，它们的涡管强度总和为$\Gamma$，它等于绕无限长机翼外不可缩周线的环量.图4.6表示双连通域中有环量流动的例子.

根据以上讨论可知，单连通域中理想不可压缩流体无旋流动只要有物面和无穷远处的流动条件，就可以求解流场；多连通域中，除了以上条件外，还需要绕不可缩周线上的环量.数学物理方程中已经证明，流体力学中上述边界条件的提法不仅是正确的，而且解是唯一的.关于拉普拉斯方程边值问题解的唯一性问题的细节，请参阅数理方程的书籍.

\section*{4.5  不可压缩无旋流动速度势方程的基本解叠加法}\addcontentsline{toc}{section}{4.5 不可压缩无旋流动速度势方程的基本解叠加法}

不可压缩无旋流动速度势满足拉普拉斯方程，拉普拉斯方程的解称为调和函数.拉普拉斯方程是线性的，因此任意调和函数$\Phi_i$的线性组合$\sum C_i\Phi_i$也是不可压缩流场速度势.如果线性组合$\sum C_i\Phi_i$的速度势又满足给定边界条件，那么这一线性叠加的调和函数$\sum C_i\Phi_i$就是给定边值问题的唯一解.利用调和函数叠加法求解不可压缩无旋流动速度势方程是一种常用的方法，特别是利用电子计算机进行不可压缩无旋流动数值求解时，从基本解叠加法发展出来的有限基本解或面元法等是目前常用的工程方法.首先介绍不可压缩无旋流场速度势方程的若干基本解.

\subsection*{1. 不可压缩流场速度势的基本解}

\textbf{(1)均匀流场}

全场速度是常数的流场称为均匀流场，也就是速度势梯度是常向量的流动：
\[
\nabla\Phi=\boldsymbol{U}_0
\]
或
\[
\frac{\partial\Phi}{\partial x}=U_0\,,\qquad \frac{\partial\Phi}{\partial y}=V_0\,,\qquad \frac{\partial\Phi}{\partial z}=W_0
\]
它满足$\Delta\Phi=0$，用全微分积分法，很容易求出它的速度势为
\begin{equation*}
\Phi=U_0\,x+V_0\,y+W_0\,z+C
\tag{4.39}
\end{equation*}
如果均匀流的方向选为$x$轴正方向，则速度势为
\[
\Phi=U_0\,x+C
\]

\textbf{(2)点源}

设在流场某一点不断有流体注入流场，其体积流量为$Q$，则这种流场称为点源流场，$Q$称作点源强度.

把坐标原点设在点源处，由于对称性，点源产生的流场只有径向速度，即
\[
\boldsymbol{U}=U(R)\,\boldsymbol{e}_R
\]
$R$为球坐标向径.围绕点源作球面，由于不可压缩流场的速度散度等于零：$\nabla\cdot\boldsymbol{U}=0$，通过以点源为球心、半径为$R$的球面的体积流量必等于点源强度，于是有
\[
4\pi U(R)\,R^{2}=Q
\]
因此不可压缩点源流动的速度场为
\[
U(R)=\frac{Q}{4\pi R^{2}}
\]
不难验证，该流场是无旋的（请读者自己证明），其速度势由速度场积分求出：
\begin{equation*}
\Phi(R)=\int\boldsymbol{U}\cdot\mathrm{d}\boldsymbol{x}=\int U(R)\,\boldsymbol{e}_R\cdot\mathrm{d}\boldsymbol{x}=\int^r U(R)\,\mathrm{d}R=-\frac{Q}{4\pi R}+C
\tag{4.40}
\end{equation*}
常数$C$不影响流场，通常可以略掉.点源解在直角坐标系中的表达式为
\[
\Phi(x,y,z)=-\frac{Q}{4\pi\sqrt{x^{2}+y^{2}+z^{2}}}
\]
设点源位于$\boldsymbol{x}_0\,(x_0,y_0,z_0)$，则任意点到点源的距离$R'$为
\[
R'=\sqrt{(x-x_0)^{2}+(y-y_0)^{2}+(z-z_0)^{2}}
\]
这时点源的速度势应为
\[
\Phi(\boldsymbol{x},\boldsymbol{x}_0)=-\frac{Q}{4\pi R'}=-\frac{Q}{4\pi\sqrt{(x-x_0)^{2}+(y-y_0)^{2}+(z-z_0)^{2}}}+C
\]
点源解的流场如图4.7所示.

如点源强度为负值，则表明在点源处不断有流体流入该点，这种负点源也称为点汇.回忆第2章中由无旋有源速度场的散度求速度场的例子，在原点给定散度等于$\delta$函数，它生成的流场和点源流场是相同的，点源$Q$相当于在点源处具有散度等于$Q\delta(\boldsymbol{x})$.

\textbf{(3)偶极子}

偶极子是由强度相等，符号相反的两点源构成的流场，当两个点源无限靠近并有$\displaystyle\lim_{\delta l\to 0}Q\delta l=m>0$时，称组合流场为偶极子流场.已经证明点源解是不可压缩有势流场，根据线性叠加原理，偶极子流场也是无旋有势的.

偶极子流动的速度势可用叠加法求得.设坐标原点处有一强度为$-Q$的点汇，在$x$轴正方向离点汇$\delta l$处有一强度相等的点源$Q$，则它们的速度势分别为
\[
\Phi_1=\frac{Q}{4\pi\sqrt{x^{2}+y^{2}+z^{2}}}\,,\qquad \Phi_2=\frac{-Q}{4\pi\sqrt{(x-\delta l)^{2}+y^{2}+z^{2}}}
\]
叠加的速度势等于
\begin{align*}
\Phi&=\lim_{\substack{\delta l\to 0\\ Q\delta l\to m}}\frac{Q}{4\pi}\bigg[\frac{1}{\sqrt{x^{2}+y^{2}+z^{2}}}-\frac{1}{\sqrt{(x-\delta l)^{2}+y^{2}+z^{2}}}\bigg]\\
&=\lim_{\substack{\delta l\to 0\\ Q\delta l\to m}}\frac{Q\delta l}{4\pi}\bigg(\frac{\partial}{\partial x}\,\frac{1}{\sqrt{x^{2}+y^{2}+z^{2}}}\bigg)
\end{align*}
完成上式的极限运算，可得偶极子流动的速度势为
\begin{equation*}
\Phi=\frac{-m}{4\pi}\cdot\frac{x}{(x^{2}+y^{2}+z^{2})^{3/2}}
\tag{4.41a}
\end{equation*}
式中$m=\displaystyle\lim_{\delta l\to 0}Q\delta l$称为偶极子强度，它有方向性，它的方向由点汇指向点源.

如果构成偶极子的点汇到点源连线在$y$轴或$z$轴上，即偶极子的指向为$y$轴或$z$轴，则相应的偶极子速度势为
\begin{equation*}
\Phi=\frac{-m}{4\pi}\cdot\frac{y}{(x^{2}+y^{2}+z^{2})^{3/2}}
\tag{4.41b}
\end{equation*}
\begin{equation*}
\Phi=\frac{-m}{4\pi}\cdot\frac{z}{(x^{2}+y^{2}+z^{2})^{3/2}}
\tag{4.41c}
\end{equation*}
在球坐标的子午面上，偶极子式(4.41a)的流线谱如图4.8所示.

空间指向为$\boldsymbol{l}=\alpha\boldsymbol{e}_1+\beta\boldsymbol{e}_2+\gamma\boldsymbol{e}_3$的偶极子的速度势可由指向$x$轴，$y$轴，$z$轴的偶极子速度势作线性叠加求出如下：
\begin{equation*}
\Phi=\frac{-m}{4\pi}\cdot\frac{\alpha x+\beta y+\gamma z}{(x^{2}+y^{2}+z^{2})^{3/2}}
\tag{4.41d}
\end{equation*}
位于空间任意点上的偶极子的速度势可通过坐标移动变换求得
\begin{equation*}
\Phi=\frac{-m}{4\pi}\cdot\frac{\alpha(x-x_0)+\beta(y-y_0)+\gamma(z-z_0)}{\big[(x-x_0)^{2}+(y-y_0)^{2}+(z-z_0)^{2}\big]^{3/2}}
\tag{4.41e}
\end{equation*}

\textbf{(4)高阶基本解及其性质}

可以通过对点源解求导的方法得到各阶基本解.当$\Phi$是不可压缩无旋流的速度势时，它的各阶偏导数
\begin{equation*}
\Phi_{lmn}=\frac{\partial^{l+m+n}\Phi}{\partial x_i^{\,l}\,\partial x_j^{\,m}\,\partial x_k^{\,n}}\,,\qquad l,m,n>0
\tag{4.42}
\end{equation*}
也满足速度势方程.事实上偶极子就是点源解的一阶偏导数.

根据点源解的特性，可以得到各阶基本解的性质如下.

\textcircled{1} 由点源解求导构成的各阶基本解在点源处具有奇异性（$\Phi\to\infty$），故又称它们为各阶奇点.在奇点处，点源解有一阶奇性，即$O(1/R)$，偶极子有二阶奇性，即$O(1/R^{2})$，$n$阶基本解有$n$阶奇性，即$O(1/R^{n})$.

\textcircled{2} 无穷远处基本解的速度势和速度渐近地趋于零，即：$\Phi_\infty=0$和$|\nabla\Phi|_\infty=0$.

\textcircled{3} 在无界域中，如将速度势在无穷远处做渐近展开，则可将其展开式写成：
\[
\Phi=D+\frac{C}{R}+C_i\,\frac{\partial}{\partial x_i}\Big(\frac{1}{R}\Big)+C_{ij}\,\frac{\partial^{2}}{\partial x_i\,\partial x_j}\Big(\frac{1}{R}\Big)+\cdots
\]
式中$R=|\boldsymbol{x}|=(x_1^{2}+x_2^{2}+x_3^{2})^{1/2}$，其中第二项是点源解，$Q=-4\pi C$是由点源注入流场的流量，第三项是偶极子解，以下各项为高阶奇点解.

以上性质对于运用各阶基本解来构造给定边界条件的速度势解很有用.第一，基本解已满足无穷远处的渐近条件，因此用基本解构造线性叠加解时，只要求满足物面上的边界条件；第二，各阶奇点必须布置在流场外，因为真实流场中不可能有速度无穷大的奇异性.作为基本解叠加法的应用，下面求解不可压缩流体绕圆球的无旋流动.

\subsection*{2. 绕圆球的不可压缩无旋流动的流场和球面压力分布}

根据凯尔文定理，当忽略外力场时，圆球在无界静止不可压缩理想流体中运动所驱动的流场是单连通域中无旋流.当圆球作匀速直线运动时，它在动力学上等价于圆球前方以均匀来流绕过圆球的定常无旋流动.将坐标原点设在球心上，该绕流问题的速度势基本方程为
\begin{equation*}
\frac{\partial^{2}\Phi}{\partial x^{2}}+\frac{\partial^{2}\Phi}{\partial y^{2}}+\frac{\partial^{2}\Phi}{\partial z^{2}}=0
\tag{4.43a}
\end{equation*}
\begin{equation*}
\frac{1}{2}\bigg[\Big(\frac{\partial\Phi}{\partial x}\Big)^{2}+\Big(\frac{\partial\Phi}{\partial y}\Big)^{2}+\Big(\frac{\partial\Phi}{\partial z}\Big)^{2}\bigg]+\frac{p}{\rho}=\frac{p_\infty}{\rho}+\frac{U_\infty^{2}}{2}
\tag{4.43b}
\end{equation*}
边界条件如下：

在半径为$R_a$的球面上
\begin{equation*}
\Big(\frac{\partial\Phi}{\partial n}\Big)_{\sqrt{x^{2}+y^{2}+z^{2}}=R_a}=0
\tag{4.43c}
\end{equation*}
无穷远处
\begin{equation*}
|\boldsymbol{x}|\to\infty:\ \frac{\partial\Phi}{\partial x}=U_\infty\,,\qquad \frac{\partial\Phi}{\partial y}=\frac{\partial\Phi}{\partial z}=0
\tag{4.43d}
\end{equation*}
基本方程(4.43a)及边界条件(4.43c)和(4.43d)已构成速度势的定解问题.因此，我们首先用基本解叠加法求出速度势，然后由式(4.43b)计算流场中和球面上的压强.

用基本解的叠加来解速度势时，可由流动基本特征来考虑选择适当的基本解.当没有绕流物体时，无穷远处来流是一个均匀流场，放入圆球后均匀绕流受到扰动，因此这一问题的解应是均匀绕流和某一扰动速度势$\phi$的叠加，即
\begin{equation*}
\Phi=U_\infty\,x+\phi
\tag{4.44}
\end{equation*}
采用哪些基本解来组成问题的速度势，需要根据问题的边界条件来选择，简单物面形状的绕流问题可由低阶奇点逐步升阶来试验.单个点源不符合封闭物形绕流问题的边界条件，因为如在圆球内有单个点源，那么球面上必有流体流出，它不符合物面上流体不可穿透条件.因此，第一个可能的选择是偶极子.根据边界条件对$x$轴的对称性，应将偶极子的轴设在$x$方向，令
\[
\phi=\frac{-mx}{4\pi\sqrt{(x^{2}+y^{2}+z^{2})^{3}}}
\]
即采用均匀流和偶极子的叠加作为绕圆球无旋流动速度势的试验：
\begin{equation*}
\Phi=U_\infty\,x-\frac{mx}{4\pi\sqrt{(x^{2}+y^{2}+z^{2})^{3}}}
\tag{4.45}
\end{equation*}
偶极子在无穷远处的速度等于零，因此该速度势满足无穷远处的来流条件，只要满足球面条件，它就是该绕流问题的唯一解.

所设的坐标系中，物面条件为：$(\partial\Phi/\partial n)_{R=R_a}=(\partial\Phi/\partial R)_{R=R_a}=0$，由式(4.45)求导数可得
\[
\frac{\partial\Phi}{\partial R}=U_\infty\,\frac{\partial x}{\partial R}-\frac{m}{4\pi}\,\frac{\partial}{\partial R}\Big(\frac{x}{R^{3}}\Big)=\Big(\frac{U_\infty}{R}+\frac{m}{2\pi R^{4}}\Big)x
\]
要求$(\partial\Phi/\partial R)_{R=R_a}=0$，得
\[
\frac{m}{4\pi}=-\frac{U_\infty R_a^{3}}{2}
\]
于是，用均匀流和强度为$m=-2\pi U_\infty R_a^{3}$的偶极子叠加的速度势满足绕圆球无旋流动的全部边界条件，因此它构成绕圆球的不可压缩无旋流的速度势解：
\begin{equation*}
\Phi=U_\infty\,x+\frac{U_\infty R_a^{3}}{2R^{3}}\,x=U_\infty\,x\Big(1+\frac{R_a^{3}}{2R^{3}}\Big)
\tag{4.46}
\end{equation*}
也可将该速度势用球坐标表达为
\begin{equation*}
\Phi=U_\infty R\cos\theta\Big(1+\frac{R_a^{3}}{2R^{3}}\Big)
\tag{4.47}
\end{equation*}
$\theta$是球坐标向径与$x$轴交角.由速度势可计算出绕圆球无旋流动的速度场为
\begin{equation*}
U_R=\frac{\partial\Phi}{\partial R}=U_\infty\cos\theta\Big(1-\frac{R_a^{3}}{R^{3}}\Big)
\tag{4.48a}
\end{equation*}
\begin{equation*}
U_\theta=\frac{1}{R}\,\frac{\partial\Phi}{\partial\theta}=-U_\infty\sin\theta\Big(1+\frac{R_a^{3}}{2R^{3}}\Big)
\tag{4.48b}
\end{equation*}
由得到的速度场可以作出不可压缩流体绕圆球无旋流动的流线谱，如图4.9所示.

球面上（$R=R_a$）流体速度为
\begin{equation*}
U_R=0\,,\qquad U_\theta=-\frac{3}{2}\,U_\infty\sin\theta
\tag{4.49a}
\end{equation*}
求出速度场后，由伯努利公式(4.43b)计算球面压强分布为
\begin{equation*}
p_{R=R_a}=p_\infty+\frac{\rho U_\infty^{2}}{2}-\rho\Big(\frac{U_R^{2}}{2}+\frac{U_\theta^{2}}{2}\Big)\bigg|_{R=R_a}=p_\infty+\frac{\rho U_\infty^{2}}{2}\Big(1-\frac{9}{4}\sin^{2}\theta\Big)
\tag{4.49b}
\end{equation*}
现在考察绕圆球不可压缩无旋流动的速度场和球面压强分布情况.

\textbf{(1)定常绕流的驻点和驻点压强}

\noindent\textbf{定义4.5}\quad 定常绕流中速度等于零的点称为驻点.在驻点处的压强称为驻点压强，驻点压强常用$p_0$表示.

由式(4.49a)可以得到：绕流物面上$\theta=0$和$\theta=\pi$处，
\[
U_R=U_\theta=0
\]
这两点称为驻点.球面上正对来流（又称迎风面）的驻点称为前驻点（$\theta=0$）；球面后部（又称背风面）的驻点称作后驻点（$\theta=\pi$）.

由式(4.49b)可以看到，定常绕流物面上驻点压强是流场中的最大压强，且
\begin{equation*}
p_0=p_\infty+\frac{\rho}{2}\,U_\infty^{2}
\tag{4.50}
\end{equation*}

\textbf{(2)最大速度和最小压强点}

由球面上速度和压强分布式(4.49a)和式(4.49b)可以断定，速度绝对值最大和压强最小值在$\theta=\pi/2$处，并且
\[
|\boldsymbol{U}|=\frac{3}{2}\,U_\infty\,,\qquad p_{\min}=p_\infty-\frac{5}{8}\,\rho U_\infty^{2}
\]

\textbf{(3)压强系数}

压强系数定义：
\[
c_p=\frac{p-p_\infty}{\dfrac{1}{2}\,\rho U_\infty^{2}}
\]
图4.10给出了绕圆球定常不可压缩无旋流动的球面压力分布和实际绕流压力分布的对比.可以看到理想的无旋绕流中，在圆球的前驻点（$\theta=0^{\circ}$）和后驻点（$\theta=180^{\circ}$）上压强最大，并等于驻点压强.圆球表面的压强分布相对于$x$轴和$y$轴都是对称的，因此作用在圆球上的压强合力等于零.和实际情况相比，在圆球迎风面上理想的压强分布和实际情况符合得相当好，而从$\theta\approx 90^{\circ}$处开始，实际压强分布开始偏离理想情况，特别是圆球后部，实际压强远低于理论压强.产生这种现象的根本原因是实际流体有粘性，因此在球面附近的流体质点不能像理想的流动那样始终贴着物面流过，而是大约在$\theta=90^{\circ}$附近流动发生分离（具体位置和流体的粘度、速度以及圆球的大小有关）.由于实际圆球绕流的分离现象，圆球表面上前后压强差的合力产生作用在圆球上的阻力（和来流速度同方向的压强合力），这种阻力通常称作压差阻力，它有别于物体表面的摩擦阻力.在第10章中将详细讨论这种由于流体粘性产生的摩擦阻力和流动分离现象.这里要强调说明的是：\textbf{在流动不发生分离或在分离点前，理想无旋绕流是真实流动的良好近似}.

基本解叠加法是求解不可压缩无旋流动速度势方程的有效方法，它的原理简单明了.根据这一原理已发展成有限基本解法以及面元法等数值计算法，在空气动力学专著中有详细介绍，本书从略.

\section*{4.6  物体在不可压缩理想流体中运动时的附加惯性}\addcontentsline{toc}{section}{4.6 物体在不可压缩理想流体中运动时的附加惯性}

本章前面例子都给出了如下结果：不可压缩理想流体均匀来流绕物体作无旋运动时，物体不承受阻力.应用惯性坐标系中动力学等价原理可以推断：物体在无界静止的不可压缩理想流体中作匀速直线运动时也不承受阻力.下面研究物体在不可压缩理想流体中作任意变速或旋转运动时受到的流体作用力.假定：

（1）不可压缩理想流体在势力场作用下由静止启动，因而流体运动是无旋的；

（2）运动物体是刚体，它只有六个运动自由度，即移动速度$\boldsymbol{U}_0$与相对于某参考点的角速度$\boldsymbol{\omega}$.

如同前面求解绕流问题一样，求解运动物体上的流体作用力问题主要是求解速度势.但是，与均匀来流的绕流问题不同，由于物体在流体中作任意运动，物体运动产生的流场速度势是非定常的，必须求出流体绝对运动的速度势.一般来说，在固定的坐标系中求解非定常速度势的初值和边值问题是比较困难的.然而，不可压缩流体无旋运动具有对边界条件的瞬时响应特性（参见4.4节中速度势的物理意义），就是说，某一瞬时的速度势只和该瞬时的边界条件有关.于是，引入\textbf{瞬时绝对坐标}的概念，用这种坐标系求解不可压缩无旋非定常流动的速度势，较之采用固定的绝对坐标系简便得多.

\subsection*{1. 瞬时绝对坐标中速度势的解}

\textbf{(1)数学提法}

由不可压缩流体在无界域中无旋流动的基本方程和边界条件(4.43a)，(4.43b)，(4.43c)和(4.43d)的表达式，可以清楚地看到，速度势$\Phi$完全由运动学边界条件确定.具体来说，在非定常情况下，某一时刻的速度势只取决于该时刻的运动边界条件，和该流场的过去历史无关.不可压缩无旋流动的这一特点，可以确切地称为流场对边界条件的瞬时响应特性.下面列出物体在不可压缩流体中作一般空间运动时的流场速度势方程和它的边界条件

速度势方程\qquad $\Delta\Phi=0$

无界静止条件\qquad $(\nabla\Phi)_\infty=\boldsymbol{0}$

物面不可穿透条件\qquad $(\nabla\Phi)_\Sigma\cdot\boldsymbol{n}=(\boldsymbol{U}_0+\boldsymbol{\omega}\times\boldsymbol{r})\cdot\boldsymbol{n}$

式中，$\boldsymbol{U}_0$是物体质心运动速度，$\boldsymbol{\omega}$是相对于质心的物体角速度，$\Sigma$是物面方程，$\boldsymbol{n}$是物面的法向量，它由物面方程$\Sigma$导出.一般情况下，以上四个量都是时间的函数.但是，由上面列出的速度势方程和边界条件可以看出：某一时刻$t_0$的速度势解只取决于该瞬时的$\boldsymbol{U}_0(t_0)$，$\boldsymbol{\omega}(t_0)$，$\Sigma(t_0)$，而和物体以前的运动状态无关.这就是说，只需要一个瞬时一个瞬时地求解速度势.利用这一特性，在每一瞬间，把坐标系冻结在物体上，并在该瞬时求出满足边界条件的绝对运动速度势，故称该坐标系为瞬时绝对坐标系.下面介绍具体解法.设瞬时冻结在刚性物体上的坐标系为$(x',y',z')$，在该坐标系中刚体物面方程和时间无关：
\begin{equation*}
\Sigma:F(x',y',z')=0
\tag{4.51}
\end{equation*}
在该坐标系中，梯度算符的表达式为
\[
\nabla'=\boldsymbol{e}'_x\,\frac{\partial}{\partial x'}+\boldsymbol{e}'_y\,\frac{\partial}{\partial y'}+\boldsymbol{e}'_z\,\frac{\partial}{\partial z'}
\]
拉普拉斯算符是
\[
\Delta'=\frac{\partial^{2}}{\partial x'^{2}}+\frac{\partial^{2}}{\partial y'^{2}}+\frac{\partial^{2}}{\partial z'^{2}}
\]
因此流动的基本方程和边界条件如下：
\begin{equation*}
\Delta'\Phi=0\,,\qquad (\nabla'\Phi)_\infty=\boldsymbol{0}\,,\qquad (\nabla'\Phi\cdot\boldsymbol{n})_\Sigma=(\boldsymbol{U}_0+\boldsymbol{\omega}\times\boldsymbol{r}')\cdot\boldsymbol{n}
\tag{4.52}
\end{equation*}
式中$\boldsymbol{n}$是物面法线：$\boldsymbol{n}=\nabla'F/|\nabla'F|$.

现在可以看到，在冻结于物体上的坐标系中，物面方程和时间无关，因此法向量$\boldsymbol{n}$也和时间无关.速度势方程的边界条件只依赖于物面方程(4.51)和物体的移动速度与转动角速度.当物形给定后，速度势方程的解仅仅线性地依赖于移动速度和角速度分量.

\textbf{(2)求解方法}

由于速度势方程是线性的，边界条件又线性地依赖于运动参数$\boldsymbol{U}_0$和$\boldsymbol{\omega}$，因此可以用以下的叠加法求解：
\begin{equation*}
\Phi=U_0\Phi_1+V_0\Phi_2+W_0\Phi_3+\omega_x\Phi_4+\omega_y\Phi_5+\omega_z\Phi_6=\sum u_i\Phi_i
\tag{4.53}
\end{equation*}
其中$(U_0,V_0,W_0)=\boldsymbol{U}_0$是物体的三个平移速度分量，$(\omega_x,\omega_y,\omega_z)=\boldsymbol{\omega}$是物体的三个角速度分量，$u_i$是广义的速度向量，它是包括平移速度分量和角速度分量的一行向量：
\[
u_i=(U_0,V_0,W_0,\omega_x,\omega_y,\omega_z)
\]
把物面边界条件展开：
\begin{align*}
\nabla'\Phi\cdot\boldsymbol{n}=\frac{\partial\Phi}{\partial n}&=\sum u_i\,\frac{\partial\Phi_i}{\partial n}=(\boldsymbol{U}_0+\boldsymbol{\omega}\times\boldsymbol{r}')\cdot\boldsymbol{n}\\
&=U_0 n_x+V_0 n_y+W_0 n_z+\omega_x\,(y'n_z-z'n_y)\\
&\quad+\omega_y\,(z'n_x-x'n_z)+\omega_z\,(x'n_y-y'n_x)
\end{align*}
可以清楚地看到，它是$u_i=(U_0,V_0,W_0,\omega_x,\omega_y,\omega_z)$的线性函数.所以，在线性叠加解(4.53)中要求每个$\Phi_i$在域内满足拉普拉斯方程；在无穷远处满足渐近条件：
\begin{equation*}
\Delta'\Phi_i=0
\tag{4.54a}
\end{equation*}
\begin{equation*}
(\nabla'\Phi_i)_\infty=\boldsymbol{0}
\tag{4.54b}
\end{equation*}
在物面上它们又分别满足：
\begin{equation*}
\begin{cases}
\dfrac{\partial\Phi_1}{\partial n}=n_x\,,\quad \dfrac{\partial\Phi_2}{\partial n}=n_y\,,\quad \dfrac{\partial\Phi_3}{\partial n}=n_z\,,\quad \dfrac{\partial\Phi_4}{\partial n}=y'n_z-z'n_y\,,\\[8pt]
\dfrac{\partial\Phi_5}{\partial n}=z'n_x-x'n_z\,,\quad \dfrac{\partial\Phi_6}{\partial n}=x'n_y-y'n_x
\end{cases}
\tag{4.54c}
\end{equation*}
以上六个速度势解的线性叠加$\sum u_i\Phi_i$就是物体一般运动的流场速度势解.从物理上来看$\Phi_i$是物体作某种单一运动的速度势.例如$\Phi_1$相当于$u_0=1,v_0=w_0=\omega_x=\omega_y=\omega_z=0$的解，即物体沿$x$轴方向以单位速度移动时速度势方程的解；$\Phi_4$则相当于$\omega_x=1,u_0=v_0=w_0=\omega_y=\omega_z=0$的速度势方程解，即物体绕$x$轴作刚体转动时流场的速度势解；以此类推，可以解释每一个$\Phi_i$所对应的物体运动边界条件.还可以看到，$\Phi_i$的定解条件式(4.54b)、(4.54c)只和物面的法向量和物面坐标有关，在冻结于物体上的坐标系中这些量和时间无关.也就是说$\Phi_i$只取决于物面几何形状，只要物面给定，$\Phi_i$只需要求解一次，这是应用瞬时绝对坐标的最大好处.求得速度势后，速度场$\boldsymbol{U}$为
\[
\boldsymbol{U}=\nabla'\Phi\,,\qquad \text{或}\quad U_{x'}=\frac{\partial\Phi}{\partial x'}\,,\quad U_{y'}=\frac{\partial\Phi}{\partial y'}\,,\quad U_{z'}=\frac{\partial\Phi}{\partial z'}
\]

\subsection*{2. 在瞬时绝对坐标系中压强计算公式}

求出速度场后，利用柯西-拉格朗日积分式
\[
\frac{\partial\Phi}{\partial t}+\frac{1}{2}\,|\boldsymbol{U}|^{2}+\frac{p}{\rho}+\Pi=C(t)
\]
可以求出压强分布.然而这时必须注意该积分式是在固定的绝对坐标系（例如，固结于地面）中导出的，具体地说，$(\partial\Phi/\partial t)_{x=\mathrm{const.}}$是在固定坐标系$\boldsymbol{x}$（而不是冻结在物体上的坐标系$\boldsymbol{x}'$）中对时间的偏导数.为此需要由$\Phi(\boldsymbol{x}',t)$求出$(\partial\Phi/\partial t)_{\boldsymbol{x}=\mathrm{const.}}$.下面通过坐标变换来导出该公式.

每一瞬时都有固定坐标系$(\boldsymbol{x},t)$和冻结坐标$(\boldsymbol{x}',t)$之间的变换关系：
\begin{equation*}
\mathrm{d}\boldsymbol{x}=\mathrm{d}\boldsymbol{x}'+\boldsymbol{U}_0\,\mathrm{d}t+\boldsymbol{\omega}\times\boldsymbol{r}'\,\mathrm{d}t
\tag{4.55}
\end{equation*}
积分该式可得到一般的坐标转换关系为
\[
\boldsymbol{x}'=\boldsymbol{x}'(\boldsymbol{x},t)
\]
或
\[
x'=x'(x,y,z,t)\,,\qquad y'=y'(x,y,z,t)\,,\qquad z'=z'(x,y,z,t)
\]
将以上坐标变换式代入冻结坐标系中的速度势表达式$\Phi'(\boldsymbol{x}',t)$（为了明确区别起见，在瞬时绝对坐标中的速度势用$\Phi'(\boldsymbol{x}',t)$表示），就可得在固定坐标中的速度势$\Phi(\boldsymbol{x},t)$：
\[
\Phi(\boldsymbol{x},t)=\Phi'\big[\boldsymbol{x}'(\boldsymbol{x},t),t\big]
\]
应用链式求导法则，可以得到
\[
\Big(\frac{\partial\Phi}{\partial t}\Big)_{\boldsymbol{x}}=\nabla'\Phi'\cdot\Big(\frac{\partial\boldsymbol{x}'}{\partial t}\Big)_{\boldsymbol{x}}+\Big(\frac{\partial'\Phi'}{\partial t}\Big)_{\boldsymbol{x}'}
\]
因$\nabla'\Phi'=\boldsymbol{U}$，而$(\partial\boldsymbol{x}'/\partial t)_{\boldsymbol{x}}$可由变换式(4.55)求出（令$\mathrm{d}\boldsymbol{x}=0$）：
\[
\Big(\frac{\partial\boldsymbol{x}'}{\partial t}\Big)_{\boldsymbol{x}}=-\boldsymbol{U}_0-\boldsymbol{\omega}\times\boldsymbol{r}'=-\boldsymbol{U}_e
\]
$\boldsymbol{U}_e$是瞬时绝对坐标系相对于固定坐标系的速度，即牵连速度，于是有
\begin{equation*}
\frac{\partial\Phi}{\partial t}=\frac{\partial'\Phi'}{\partial t}-\boldsymbol{U}\cdot\boldsymbol{U}_e=\frac{\partial'\Phi'}{\partial t}-\nabla'\Phi'\cdot\boldsymbol{U}_e
\tag{4.56}
\end{equation*}
式中$\boldsymbol{U}_e=\boldsymbol{U}_0+\boldsymbol{\omega}\times\boldsymbol{r}'$是瞬时绝对坐标系当地的牵连速度.代入柯西-拉格朗日积分，可得到瞬时绝对坐标系中的表达式：
\[
\frac{\partial'\Phi'}{\partial t}-\boldsymbol{U}_e\cdot\nabla'\Phi'+\frac{1}{2}\,|\boldsymbol{U}|^{2}+\frac{p}{\rho}+\Pi=C(t)
\]
现在公式中的函数都是用瞬时绝对坐标系表示的，因此下面将速度势符号上的撇号取消.

最后，确定柯西-拉格朗日积分式的常数$C(t)$，前面已经提到过，无界静止的无源流场中速度势在无穷远处的渐近式为
\[
\Phi\approx C_j\,\frac{\partial}{\partial x_j}\Big(\frac{1}{R}\Big)+\cdots=O\Big(\frac{1}{R^{2}}\Big)+\text{高阶量}
\]
因而$\partial\Phi/\partial t|_\infty=0$和$|\nabla\Phi|_\infty=0$.如果选定无穷远处$\Pi=0$，将以上等式代入柯西-拉格朗日积分式，则可得积分常数$C(t)$为
\[
C(t)=\frac{p_\infty}{\rho}
\]
于是在瞬时绝对坐标系中压强场的解为
\begin{equation*}
\frac{p}{\rho}=\frac{p_\infty}{\rho}-\Big(\frac{\partial\Phi}{\partial t}+\frac{U^{2}}{2}+\Pi-\boldsymbol{U}\cdot\boldsymbol{U}_e\Big)
\tag{4.57}
\end{equation*}
下面举例说明以上方法的应用.

\noindent\textbf{例4.4}\quad 圆球在不可压缩理想流体中作加速直线运动时，求作用在圆球上的合力（不计外力场）.

\noindent\textbf{解}\quad 在冻结于圆球的坐标系（图4.11）中求$\Phi_1$，基本方程和边界条件如下：
\[
\Delta'\Phi_1=0\,,\qquad (\nabla'\Phi_1)_\infty=\boldsymbol{0}\,,\qquad \Big(\frac{\partial\Phi_1}{\partial n}\Big)_{R_a}=n_x=\frac{x'}{R_a}=\cos\theta'
\]
$R_a$为圆球半径.

前面已经导出均匀来流绕圆球速度势，它等于均匀流和偶极子速度势的叠加.在冻结于圆球上的坐标系中，无穷远处速度等于零，因此在瞬时绝对坐标系中圆球平移运动产生的流场速度势$\Phi_1$就是偶极子解：
\[
\Phi_1=\frac{-R_a^{3}x'}{2R'^{3}}=-\frac{R_a^{3}\cos\theta'}{2R'^{2}}
\]
验证边界条件：该速度势在球面上的法向导数
\[
\frac{\partial\Phi_1}{\partial R'}\bigg|_{R_a}=\frac{R_a^{3}\cos\theta'}{R'^{3}}\bigg|_{R_a}=\cos\theta'
\]
它满足球面的边界条件.当圆球以速度$U_0(t)$运动时产生的速度势为
\[
\Phi=\frac{-U_0(t)\,R_a^{3}x'}{2R'^{3}}=-\frac{U_0(t)\,R_a^{3}\cos\theta'}{2R'^{2}}
\]
因而在瞬时绝对坐标系中速度势对时间的偏导数
\[
\Big(\frac{\partial\Phi}{\partial t}\Big)_{x'}=-\dot{U}_0(t)\,\frac{R_a^{3}x'}{2R'^{3}}=-\frac{\dot{U}_0(t)\,R_a^{3}}{2}\,\frac{\cos\theta'}{R'^{2}}
\]
球面上流体速度
\[
\boldsymbol{U}=(\nabla'\Phi)_{R_a}=U_0\Big(\cos\theta'\,\boldsymbol{e}_{R'}+\frac{1}{2}\sin\theta'\,\boldsymbol{e}_{\theta'}\Big)\,,\qquad U^{2}=U_0^{2}\Big(1-\frac{3}{4}\sin^{2}\theta'\Big)
\]
圆球面上的牵连速度为：$\boldsymbol{U}_e=U_0\boldsymbol{e}_{x'}$，因此
\[
\boldsymbol{U}\cdot\boldsymbol{U}_e=U_0\Big(\cos\theta'\,\boldsymbol{e}_{R'}+\frac{1}{2}\sin\theta'\,\boldsymbol{e}_{\theta'}\Big)\cdot\big(U_0\,\boldsymbol{e}_{x'}\big)=U_0^{2}\Big(1-\frac{3}{2}\sin^{2}\theta'\Big)
\]
代入柯西-拉格朗日积分式后得
\[
\frac{p}{\rho}=\frac{p_\infty}{\rho}-\frac{U_0^{2}}{2}\Big(1-\frac{9}{4}\cos^{2}\theta'\Big)+\frac{R_a}{2}\,\dot{U}_0\cos\theta'
\]
和圆球作等速直线运动的结果比较（式(4.49)），球面上的压强增加一项$\dfrac{\mathrm{d}U_0}{\mathrm{d}t}\dfrac{R_a\cos\theta'}{2}$，这是圆球加速运动时在球面上产生的流体压强，这一部分压强分布相对于$x$轴反对称，因此将产生沿速度方向的合力.压强分布中其他各项相对于$x$轴对称，并和理想不可压缩流体绕圆球无旋流动的压强分布相同，这部分压强分布的合力等于零.总合力计算如下：
\begin{align*}
\boldsymbol{F}&=-\iint p\boldsymbol{n}\,\mathrm{d}A=-\iint\big(p\cos\theta'\,\boldsymbol{e}_x+p\sin\theta'\,\boldsymbol{e}_y\big)\mathrm{d}A\\
&=-2\pi\int_0^{\pi}\bigg[\frac{p_\infty}{\rho}-\frac{U_0^{2}}{2}\Big(1-\frac{9}{4}\sin^{2}\theta'\Big)+\frac{1}{2}\,\rho\,\frac{\mathrm{d}U_0}{\mathrm{d}t}\,R_a\cos\theta'\bigg]\big(\cos\theta'\,\boldsymbol{e}_x+\sin\theta'\,\boldsymbol{e}_y\big)R_a^{2}\sin\theta'\,\mathrm{d}\theta'\\
&=-\frac{2}{3}\,\pi\rho R_a^{3}\,\frac{\mathrm{d}U_0}{\mathrm{d}t}\,\boldsymbol{e}_x
\end{align*}
该合力方向和圆球运动方向相反，因此是阻力.阻力的大小正比于圆球的加速度、流体密度和圆球体积之半（$2\pi R_a^{3}/3$）.也就是说该阻力相当于在圆球中附加了半球体积的流体质量（$2\pi\rho R_a^{3}/3$）后作同样运动所需的惯性力，因而$2\pi\rho R_a^{3}/3$称为圆球体在流体中运动的附加质量.

该例题表明物体在流体中作直线加速度运动时，除了物体具有惯性外，流体的加速运动也表现某种惯性效应，或称附加惯性.下面介绍物体作一般运动时（移动和转动）附加惯性的概念和公式.

\subsection*{3. 物体在不可压缩理想流体中运动时的附加惯性}

应用能量定理来导出物体在不可压缩理想流体中作任意运动时受到的流体作用力和力矩.首先导出流体运动的能量守恒公式.

\textbf{(1)无界不可压缩理想流体无旋运动的能量守恒公式}

不可压缩理想流体质量体上的能量方程为
\[
\frac{D}{Dt}\int_{V(t)}\rho\,\frac{|\boldsymbol{U}|^{2}}{2}\,\mathrm{d}V=\int_{V(t)}\rho\boldsymbol{f}\cdot\boldsymbol{U}\,\mathrm{d}V+\oint_{\Sigma'(t)}\boldsymbol{T}_n\cdot\boldsymbol{U}\,\mathrm{d}A
\]
对于理想流体，流体应力状态$\boldsymbol{T}_n=-p\boldsymbol{n}$，$p$是压强，当不计外力时，理想流体运动的能量方程为
\[
\frac{D}{Dt}\int_{V(t)}\rho\,\frac{|\boldsymbol{U}|^{2}}{2}\,\mathrm{d}V=-\oint_{\Sigma'(t)}p\boldsymbol{n}\cdot\boldsymbol{U}\,\mathrm{d}A
\]
式中$\boldsymbol{n}$是$\Sigma'(t)$的外法线方向.当物体在无界静止流体中运动时，能量方程右端的压强做功项包括两部分：一部分是物体表面$\Sigma(t)$上压强做功；还有一部分是半径为无穷大的封闭曲面压强做功.前面已经证明，在无源的无旋流场中，$|\nabla\Phi|_{\boldsymbol{x}\to\infty}=0$和$\partial\Phi/\partial l|_{\boldsymbol{x}\to\infty}=0$，压强等于常数.因而，压强做功的面积分中，无穷大封闭曲面上积分值等于零.也就是说，\textbf{物体在理想不可压缩流体中运动时，全流场的动能增长率等于物体表面上压强所做的功率}.
\[
\frac{D}{Dt}\int_{V(t)}\rho\,\frac{|\boldsymbol{U}|^{2}}{2}\,\mathrm{d}V=-\oint_{\Sigma(t)}p\boldsymbol{n}\cdot\boldsymbol{U}\,\mathrm{d}A
\]
根据物体表面上的边界条件（瞬时绝对坐标的撇号略去），
\begin{align*}
\boldsymbol{U}\cdot\boldsymbol{n}&=(\boldsymbol{U}_0+\boldsymbol{\omega}\times\boldsymbol{r})\cdot\boldsymbol{n}=U_0 n_x+V_0 n_y+W_0 n_z\\
&\quad+\omega_x\,(y n_z-z n_y)+\omega_y\,(z n_x-x n_z)+\omega_z\,(x n_y-y n_x)
\end{align*}
将它代入表面力做功的面积分中，有
\begin{align*}
-\oint_{\Sigma(t)}p\boldsymbol{n}\cdot\boldsymbol{U}\,\mathrm{d}A&=U_0\oint_{\Sigma(t)}-p n_x\,\mathrm{d}A+V_0\oint_{\Sigma(t)}-p n_y\,\mathrm{d}A+W_0\oint_{\Sigma(t)}-p n_z\,\mathrm{d}A\\
&\quad+\omega_x\oint_{\Sigma(t)}-p\,(y n_z-z n_y)\,\mathrm{d}A+\omega_y\oint_{\Sigma(t)}-p\,(z n_x-x n_z)\,\mathrm{d}A\\
&\quad+\omega_z\oint_{\Sigma(t)}-p\,(x n_y-y n_x)\,\mathrm{d}A
\end{align*}
我们知道面积分$\displaystyle\oint_{\Sigma(t)}-p n_x\,\mathrm{d}A$是物体作用在流体上合力的$x$方向分量；其他类似的二项分别是合力的$y$方向和$z$方向的分量；$\displaystyle\oint_{\Sigma(t)}-p\,(y n_z-z n_y)\,\mathrm{d}A$等于物体作用于流体上的合力矩的$x$方向分量，类似地，其他二项分别是合力矩的$y$方向和$z$方向的分量.如用$\boldsymbol{F}=F_x\boldsymbol{i}+F_y\boldsymbol{j}+F_z\boldsymbol{k}$表示合力；$\boldsymbol{L}=L_x\boldsymbol{i}+L_y\boldsymbol{j}+L_z\boldsymbol{k}$表示合力矩，则有
\[
F_x=\oint_{\Sigma(t)}-p n_x\,\mathrm{d}A\,,\qquad F_y=\oint_{\Sigma(t)}-p n_y\,\mathrm{d}A\,,\qquad F_z=\oint_{\Sigma(t)}-p n_z\,\mathrm{d}A
\]
\[
L_x=\oint_{\Sigma(t)}-p\,(y n_z-z n_y)\,\mathrm{d}A\,,\qquad L_y=\oint_{\Sigma(t)}-p\,(z n_x-x n_z)\,\mathrm{d}A
\]
\[
L_z=\oint_{\Sigma(t)}-p\,(x n_y-y n_x)\,\mathrm{d}A
\]
代入能量积分公式，可得
\[
\frac{D}{Dt}\int_{V(t)}\rho\,\frac{|\boldsymbol{U}|^{2}}{2}\,\mathrm{d}V=F_x U_0+F_y V_0+F_z W_0+L_x\omega_x+L_y\omega_y+L_z\omega_z
\]
上面导出的能量守恒公式可表述为：不计质量力时

\textbf{作用在理想流体质量体上的外界功率等于流场总动能的增长率}

下面计算流场的总动能，然后代入能量积分公式，就可以求出在物体上的流体作用力.

\textbf{(2)理想不可压缩流体无旋运动的全场动能公式}

不可压缩流体的无界无旋流动的流场总动能
\[
T=\frac{\rho}{2}\int\boldsymbol{U}\cdot\boldsymbol{U}\,\mathrm{d}V=\frac{\rho}{2}\iiint(\nabla\Phi\cdot\nabla\Phi)\,\mathrm{d}V
\]
应用向量导数公式：$\nabla\Phi\cdot\nabla\Phi=\nabla\cdot(\Phi\nabla\Phi)-\Phi\Delta\Phi$，因为不可压缩无旋流的速度势满足拉普拉斯方程：$\Delta\Phi=0$，故$\nabla\Phi\cdot\nabla\Phi=\nabla\cdot(\Phi\nabla\Phi)$.利用高斯公式，动能公式可简化为
\[
T=\frac{\rho}{2}\iiint\nabla\Phi\cdot\nabla\Phi\,\mathrm{d}V=\frac{\rho}{2}\iiint\nabla\cdot(\Phi\nabla\Phi)\,\mathrm{d}V=\frac{\rho}{2}\oint_\Sigma\Phi\,\frac{\partial\Phi}{\partial n'}\,\mathrm{d}A+\frac{\rho}{2}\oint_S\Phi\,\frac{\partial\Phi}{\partial r}\,\mathrm{d}A
\]
式中$\Sigma$为物面，$\boldsymbol{n}'$是指向物面的内法线方向（流场的外法方向）.$S$是半径足够大的圆球面，由于不可压缩无源无旋流场中$\displaystyle\lim_{r\to\infty}\Phi\sim O\Big(\frac{1}{r^{2}}\Big)$，$\displaystyle\lim_{r\to\infty}\frac{\partial\Phi}{\partial r}\sim O\Big(\frac{1}{r^{3}}\Big)$，因此当圆球面取得无限大时，第二个积分等于零.此外，取指向物面外的法线方向$\boldsymbol{n}=-\boldsymbol{n}'$，则动能公式的最后形式为
\begin{equation*}
T=-\frac{\rho}{2}\oint_\Sigma\Phi\,\frac{\partial\Phi}{\partial n}\,\mathrm{d}A
\tag{4.58}
\end{equation*}
将物体一般运动产生的速度势公式$\Phi=\sum u_i\Phi_i$代入，得
\[
T=-\frac{\rho}{2}\sum_{i=1}^{6}\sum_{j=1}^{6}u_i u_j\oint_\Sigma\Phi_i\,\frac{\partial\Phi_j}{\partial n}\,\mathrm{d}A
\]
令
\begin{equation*}
a_{ij}=-\oint_\Sigma\Phi_i\,\frac{\partial\Phi_j}{\partial n}\,\mathrm{d}A
\tag{4.59}
\end{equation*}
则
\begin{equation*}
T=\frac{\rho}{2}\sum_{i=1}^{6}\sum_{j=1}^{6}a_{ij}\,u_i u_j
\tag{4.60}
\end{equation*}
式(4.60)是不可压缩无旋流场动能的一般公式.前面已经论述过$\Phi_i$只与物面形状有关，因此$a_{ij}$只是物面形状的函数而与物体运动参数无关，称$a_{ij}$为流体的附加惯性张量，而且有以下主要性质：

\textcircled{1} $a_{ij}$是正定矩阵，因为动能总是正值.

\textcircled{2} $a_{ij}=a_{ji}$是对称矩阵.

在以半径无穷大的球面包围的流场中应用格林公式，有
\begin{align*}
\iiint_V\Big(\Phi_i\Delta\Phi_j-\Phi_j\Delta\Phi_i\Big)\mathrm{d}V&=\iiint_V\big[\nabla\cdot(\Phi_i\nabla\Phi_j)-\nabla\cdot(\Phi_j\nabla\Phi_i)\big]\mathrm{d}V\\
&=\oint_\Sigma\Big(\Phi_i\,\frac{\partial\Phi_j}{\partial n}-\Phi_j\,\frac{\partial\Phi_i}{\partial n}\Big)\mathrm{d}A+\oint_{S(r\to\infty)}\Big(\Phi_i\,\frac{\partial\Phi_j}{\partial r}-\Phi_j\,\frac{\partial\Phi_i}{\partial r}\Big)\mathrm{d}A
\end{align*}
由于$\displaystyle\lim_{r\to\infty}\Phi_i\sim O\Big(\frac{1}{r^{2}}\Big)$和$\displaystyle\lim_{r\to\infty}\frac{\partial\Phi}{\partial r}\sim O\Big(\frac{1}{r^{3}}\Big)$，故第二个面积分等于零，同时$\Delta\Phi_i=\Delta\Phi_j=0$，故有
\[
a_{ij}=-\oint_\Sigma\Phi_i\,\frac{\partial\Phi_j}{\partial n}\,\mathrm{d}A=-\oint_\Sigma\Phi_j\,\frac{\partial\Phi_i}{\partial n}\,\mathrm{d}A=a_{ji}
\]
证毕.

由于$a_{ij}$的对称性，事实上$a_{ij}$只有21个独立分量.

\textcircled{3} 若物面具有三个相互垂直的对称面（例如椭球：$x^{2}/a^{2}+y^{2}/b^{2}+z^{2}/c^{2}=1$），则该物体在不可压缩理想流体中作无旋运动时
\[
a_{ij}=-\rho\oint_\Sigma\Phi_i\,\frac{\partial\Phi_j}{\partial n}\,\mathrm{d}A=0\,,\qquad i\neq j
\]
设对称物面的方程为$F(x,y,z)=0$，由于对称性，应有$F(x,y,z)=F(-x,y,z)$，$F(x,y,z)=F(x,-y,z)$以及$F(x,y,z)=F(x,y,-z)$等；但是物面的法向量的分量有反对称性，例如$n_y(x,y,z)=\partial F/\partial y\,/\,|\nabla F|=-n_y(x,-y,z)$.

现在以$a_{12}=-\rho\displaystyle\oint_\Sigma\Phi_1(x,y,z)\dfrac{\partial\Phi_2}{\partial n}\,\mathrm{d}A$为例，证明$a_{12}=0$.物体以单位速度沿$x$轴正方向运动时的流场速度势为$\Phi_1=\Phi_1(x,y,z)$，同时根据定义，物面上$\partial\Phi_2/\partial n=n_y$，它是物面单位法向量在$y$轴方向的分量，因此有
\[
a_{12}=-\rho\oint_\Sigma n_y\,\Phi_1(x,y,z)\,\mathrm{d}A
\]
取对称点$M(x,y,z)$和$M'(x,-y,z)$（参见图4.12），因为物面对$xz$平面对称，故物体沿$x$轴运动的速度势$\Phi_1=\Phi_1(x,y,z)$也对$xz$平面对称，即$\Phi_1(x,y,z)=\Phi_1(x,-y,z)$.另一方面，对称物面的法向量分量有反对称性：$n_y(x,y,z)=(\partial F/\partial y)/|\nabla F|=-n_y(x,-y,z)$，于是，在面积分式$a_{12}=-\rho\displaystyle\oint_\Sigma\Phi_1\dfrac{\partial\Phi_2}{\partial n}\,\mathrm{d}A$中，在$xz$平面上方的面积分值和$xz$平面下方的面积分值相互抵消，即$a_{12}=0$.同理可证其他的交叉惯性张量分量$a_{ij}=0\,(i\neq j)$.

于是有三个相互垂直对称面的物体$a_{ij}$只有六个独立量，总动能可写成
\[
T=\frac{\rho}{2}\big(a_{11}U_0^{2}+a_{22}V_0^{2}+a_{33}W_0^{2}+a_{44}\omega_x^{2}+a_{55}\omega_y^{2}+a_{66}\omega_z^{2}\big)
\]

\textbf{(3)不可压缩理想流体无界无旋流场中物体上的合力、合力矩和达朗贝尔定理}

现在来考察物体在不可压缩理想流体中任意运动时的受力和运动.将动能公式代入动能守恒公式，得
\begin{equation*}
f_i u_i=\frac{\rho}{2}\,\frac{\mathrm{d}}{\mathrm{d}t}\sum a_{ij}\,u_i u_j
\tag{4.61}
\end{equation*}
式中$f_i=(F_x,F_y,F_z,L_x,L_y,L_z)$是一行力向量；$u_i=(U_0,V_0,W_0,\omega_x,\omega_y,\omega_z)$是一行速度向量.将求导数与求和运算交换，由于$a_{ij}=a_{ji}$，并和时间无关，式(4.61)可简化为
\begin{equation*}
f_i u_i=\rho\sum a_{ij}\,u_i\,\frac{\mathrm{d}u_j}{\mathrm{d}t}
\tag{4.62}
\end{equation*}
上式对于任意$u_i$都成立，因而有
\begin{equation*}
f_i=\rho\sum a_{ij}\,\frac{\mathrm{d}u_j}{\mathrm{d}t}
\tag{4.63}
\end{equation*}
这是流体上受到物体作用力和力矩的一般公式.根据力的作用和反作用原理，流体作用在物体上的力和力矩等于$-f_i$.

由式(4.63)可以很容易导出著名的达朗贝尔定理：\textbf{物体在理想不可压缩流体中作匀速平移运动时不受阻力}.

当物体作等速平移运动时，$\boldsymbol{\omega}=\boldsymbol{0}$，$\boldsymbol{U}_0=U\boldsymbol{e}_x$为常向量，因此$\mathrm{d}u_i/\mathrm{d}t=0$.则由式(4.63)可得
\[
F_x=\rho\sum a_{xx}\,\frac{\mathrm{d}U}{\mathrm{d}t}=0
\]
就是说物体在理想不可压缩流体中作等速平移运动时，物体上不受任何作用力或力矩.真实流体有粘性而不是理想流体，因此即使是无分离的绕流，物体上还是受到流体的摩擦阻力，关于粘性流体作用在物体上的阻力问题将在第8章和第10章中讨论.

下面进一步讨论物体在流体中作变速直线运动和转动的情况.

如果物体是有三个互相垂直的对称面，则在直角坐标中
\[
a_{ij}\equiv a_{ji}=0
\]
因此
\[
f_\alpha=\rho\,a_{\alpha\alpha}\,\frac{\mathrm{d}u_\alpha}{\mathrm{d}t}\qquad(\text{对}\,\alpha\,\text{不求和})
\]
当物体受到外加作用力$\boldsymbol{R}=(R_i)$和力矩$\boldsymbol{M}=(M_i)$推动下在流体中运动时，物体的运动方程为
\begin{equation*}
M\,\frac{\mathrm{d}U_\alpha}{\mathrm{d}t}=R_\alpha+F_\alpha
\tag{4.64}
\end{equation*}
\begin{equation*}
I_\alpha\,\frac{\mathrm{d}\omega_\alpha}{\mathrm{d}t}=M_\alpha+L_\alpha\qquad(\text{对}\,\alpha\,\text{不求和})
\tag{4.65}
\end{equation*}
$M$为物体质量，$I_\alpha$为物体相对于坐标轴的转动惯量.将流体作用力公式(4.63)代入，得
\begin{equation*}
\begin{cases}
M\,\dfrac{\mathrm{d}U_\alpha}{\mathrm{d}t}=R_\alpha-\rho a_{\alpha\alpha}\,\dfrac{\mathrm{d}U_\alpha}{\mathrm{d}t}\,,\quad \alpha=1,2,3\\[8pt]
I_\alpha\,\dfrac{\mathrm{d}\omega_\alpha}{\mathrm{d}t}=M_\alpha-\rho a_{\alpha\alpha}\,\dfrac{\mathrm{d}\omega_\alpha}{\mathrm{d}t}\,,\quad \alpha=4,5,6
\end{cases}
\qquad(\text{对}\,\alpha\,\text{不求和})
\tag{4.66}
\end{equation*}
写成分量形式
\[
(M+\rho a_{11})\,\frac{\mathrm{d}U_0}{\mathrm{d}t}=R_x\,,\qquad (M+\rho a_{22})\,\frac{\mathrm{d}V_0}{\mathrm{d}t}=R_y\,,\qquad (M+\rho a_{33})\,\frac{\mathrm{d}W_0}{\mathrm{d}t}=R_z
\]
\[
(I_x+\rho a_{44})\,\frac{\mathrm{d}\omega_x}{\mathrm{d}t}=M_x\,,\qquad (I_y+\rho a_{55})\,\frac{\mathrm{d}\omega_y}{\mathrm{d}t}=M_y\,,\qquad (I_z+\rho a_{66})\,\frac{\mathrm{d}\omega_z}{\mathrm{d}t}=M_z
\]
可以看到，$\rho a_{\alpha\alpha}$相当于在物体上增加质量或惯性矩，因此称它为附加惯性，也就是说物体在不可压缩理想流体中作加速运动时不仅要克服自身的惯性，还需要附加克服流体惯性的作用力和力矩.利用附加惯性可以使理想流体对物体加速运动时作用力的计算得到简化，对于每一种物形可以先算出$a_{ij}$（它们只与物体外形有关），然后代入物体运动方程(4.66)，就可以求解刚体运动常微分方程.

附加质量和流体密度成正比，物体在空气中运动时的附加质量远远小于水中的附加质量，因此物体在空气中运动时附加质量往往可以忽略.如果运动物面外形没有对称性，这时交叉附加惯性$a_{ij}\neq 0\,(i\neq j)$，这就意味着，在物体某一运动方向上流体作用力不仅和同方向的附加惯性有关，还和其他方向的附加惯性及加速度有关.六个刚体运动方程是耦合的，物体的移动速度可能产生流体的作用力矩而引起物体的转动.这就是说，非对称物体在流体中运动时，它的运动姿态较之在真空中复杂得多.

\noindent\textbf{例4.5}\quad 质量为$m$，直径为$d$的圆球悬挂在长为$l$的细丝上，放在密度为$\rho$的理想不可压缩流体里，并在重力场中作微幅摆动，求该单摆的摆动周期.

\noindent\textbf{解}\quad 该摆上垂直于地面的力为重力和浮力之和，即
\[
F_z=mg-\frac{\rho g\,\pi d^{3}}{6}
\]
切向加速度为$l\,\dfrac{\mathrm{d}\omega}{\mathrm{d}t}=l\,\dfrac{\mathrm{d}^{2}\theta}{\mathrm{d}t^{2}}$，平衡方程为
\[
l\,(m+\rho a_{\theta\theta})\,\frac{\mathrm{d}^{2}\theta}{\mathrm{d}t^{2}}=-\Big(mg-\frac{\rho\pi d^{3}g}{6}\Big)\theta
\]
已知圆球的$a_{\theta\theta}=\pi d^{3}/12$，得球摆的运动方程为
\[
\frac{\mathrm{d}^{2}\theta}{\mathrm{d}t^{2}}+\frac{g\big(m-\rho\pi d^{3}/6\big)\theta}{l\big(m+\rho\pi d^{3}/12\big)}=0
\]
上式可写成$\mathrm{d}^{2}\theta/\mathrm{d}t^{2}+\omega^{2}\theta=0$的形式，$\omega^{2}=g\,(m-\rho\pi d^{3}/6)\,/\,\big[l\,(m+\rho\pi d^{3}/12)\big]$，故该方程的解为周期运动$\theta=A\sin\omega t+B\cos\omega t$，周期$T=2\pi/\omega$.因此考虑流体附加惯性的单摆的周期为
\[
T=2\pi\sqrt{\frac{l\big(m+\rho\pi d^{3}/12\big)}{g\big(m-\rho\pi d^{3}/6\big)}}
\]
不计流动阻力，该摆的摆动周期比真空中的周期$\sqrt{l/g}$大$\sqrt{\dfrac{12m+\rho\pi d^{3}}{12m-2\rho\pi d^{3}}}$倍，解毕.

下面给出一些简单物形的附加惯性.

表4.1中圆柱的附加质量是轴向单位长度的附加质量；旋转椭球附加质量中系数$k_1,k_2$的值由表4.2给出.

\begin{center}
\captionof{table}{圆球、旋转椭球和圆柱的附加惯性（乘密度$\rho$）}
\begin{tabular}{|c|c|c|c|c|c|}
\hline
物体 & 几何量 & $a_{xx}$ & $a_{yy}$ & $a_{zz}$ & $a_{ij}\,(i\neq j)$\\
\hline
圆球 & 圆球半径$R_a$ & $2\pi R_a^{3}/3$ & $2\pi R_a^{3}/3$ & $2\pi R_a^{3}/3$ & $0$\\
\hline
旋转椭球 & 旋转椭球长轴$a$，旋转半径$b$ & $4\pi k_2 ab^{2}/3$ & $4\pi k_1 ab^{2}/3$ & $4\pi k_1 ab^{2}/3$ & $0$\\
\hline
圆柱 & 圆柱半径$a$ & $\pi a^{2}$ & $\pi a^{2}$ & $\infty$ & $0$\\
\hline
\end{tabular}
\end{center}

\begin{center}
\captionof{table}{旋转椭球的附加质量系数}
\begin{tabular}{|c|c|c|c|c|c|c|c|c|c|c|c|}
\hline
$a/b$ & 1.0 & 1.5 & 2.0 & 2.5 & 3.0 & 4.0 & 5.0 & 6.0 & 8.0 & 10.0 & $\infty$\\
\hline
$k_1$ & 0.5 & 0.621 & 0.702 & 0.763 & 0.803 & 0.860 & 0.895 & 0.918 & 0.945 & 0.960 & 1.0\\
\hline
$k_2$ & 0.5 & 0.304 & 0.209 & 0.156 & 0.122 & 0.082 & 0.059 & 0.045 & 0.029 & 0.021 & 0.0\\
\hline
\end{tabular}
\end{center}

最后，还应再次强调指出，本章关于流动中流体对物体作用力的分析和结论都是以理想的无粘性流体为前提.关于粘性的作用以及有关阻力等现象，将在第8章和第10章中介绍.

\section*{4.7  理想不可压缩流体中的涡动力学}\addcontentsline{toc}{section}{4.7 理想不可压缩流体中的涡动力学}

上面讨论了理想不可压缩流体的无旋流动，根据凯尔文定理，它是在有势力场作用下正压流体从无旋的初始条件下产生的.如果流体是非正压或外力场非有势，则初始的无旋流场可以发展为有旋流动.下面简要地讨论理想流体中有旋流动的产生和它的主要性质.

\subsection*{1. 理想流体的涡动力学方程}

理想流体的动力学方程为
\[
\frac{\partial\boldsymbol{U}}{\partial t}+\boldsymbol{U}\cdot\nabla\boldsymbol{U}=-\frac{1}{\rho}\,\nabla p+\boldsymbol{f}
\]
其中$\boldsymbol{f}$为外力场，注意到向量场公式：$\boldsymbol{U}\cdot\nabla\boldsymbol{U}=\nabla\,|\boldsymbol{U}|^{2}/2-\boldsymbol{U}\times(\nabla\times\boldsymbol{U})$，以及$\boldsymbol{\omega}=\nabla\times\boldsymbol{U}$，对运动方程取旋度后，有如下涡量的动力学方程：
\begin{equation*}
\frac{\partial\boldsymbol{\omega}}{\partial t}+\nabla\times(\boldsymbol{U}\times\boldsymbol{\omega})=-\nabla\times\Big(\frac{1}{\rho}\,\nabla p\Big)+\nabla\times\boldsymbol{f}
\tag{4.67}
\end{equation*}
再利用向量公式
\begin{gather*}
\nabla\times(\boldsymbol{U}\times\boldsymbol{\omega})=\boldsymbol{U}(\nabla\cdot\boldsymbol{\omega})-\boldsymbol{\omega}(\nabla\cdot\boldsymbol{U})-\boldsymbol{U}\cdot\nabla\boldsymbol{\omega}+\boldsymbol{\omega}\cdot\nabla\boldsymbol{U}\\
\nabla\times\Big(\frac{1}{\rho}\,\nabla p\Big)=\Big(\nabla\frac{1}{\rho}\Big)\times\nabla p+\frac{1}{\rho}\,\nabla\times(\nabla p)=-\frac{1}{\rho^{2}}\,\nabla\rho\times\nabla p
\end{gather*}
理想流体的涡动力学方程可写作
\begin{equation*}
\frac{\partial\boldsymbol{\omega}}{\partial t}+\boldsymbol{U}\cdot\nabla\boldsymbol{\omega}=\boldsymbol{\omega}\cdot\nabla\boldsymbol{U}-\boldsymbol{\omega}\,(\nabla\cdot\boldsymbol{U})+\frac{\nabla\rho\times\nabla p}{\rho^{2}}+\nabla\times\boldsymbol{f}
\tag{4.68}
\end{equation*}
方程(4.68)表明沿流体质点轨迹涡量的增长率由以下四部分作出贡献：①在切变场中涡量的生成项$\boldsymbol{\omega}\cdot\nabla\boldsymbol{U}$；②因微团体积膨胀使涡量减少$-\boldsymbol{\omega}\nabla\cdot\boldsymbol{U}$，对于不可压缩流体，该项等于零；③非正压流场$\nabla\rho\times\nabla p\neq\boldsymbol{0}$，对于不可压缩流体该项等于零；④外力场的旋度，对于有势力场，该项等于零.

下面分别来讨论这四部分的性质.

\subsection*{2. 涡量生成的讨论}

\textbf{(1)切变场中涡量生成项$\boldsymbol{\omega}\cdot\nabla\boldsymbol{U}$}

首先考察$\boldsymbol{\omega}\cdot\nabla\boldsymbol{U}$的性质.如果取一微元涡管作为考察的微元流体质量体（图4.13），涡管$AB$上$B$点相对于$A$有速度差$\delta\boldsymbol{U}$（图4.13(a)），对于该微元涡管
\[
\boldsymbol{\omega}\cdot\nabla\boldsymbol{U}=|\boldsymbol{\omega}|\,\delta U=|\boldsymbol{\omega}|\,\delta U_\perp+|\boldsymbol{\omega}|\,\delta U_\parallel
\]
其中$\delta\boldsymbol{U}_\perp$为$\delta\boldsymbol{U}$垂直于$\boldsymbol{\omega}$的分量，$\delta\boldsymbol{U}_\parallel$为平行于$\boldsymbol{\omega}$的分量.$\delta\boldsymbol{U}_\perp$使涡管绕$A$点转动，$\delta\boldsymbol{U}_\parallel$使涡管伸长（$\boldsymbol{\omega}\cdot\nabla\boldsymbol{U}>0$）或收缩（$\boldsymbol{\omega}\cdot\nabla\boldsymbol{U}<0$）.转动部分在当地并不改变涡量的绝对值，只改变它的方向.涡管伸长将使涡量增加，因为当微元涡管拉长时，如不计体积的变化（体积变化在另一项中考虑），截面积将按反比减小.由于涡管通量守恒性，$\boldsymbol{\omega}$的绝对值将增大；同理，涡管收缩，$\boldsymbol{\omega}$的绝对值将减小.由上面的分析可知，当涡管处于切变场时，它的方向和大小都可能发生变化，这是由涡量方程中$\boldsymbol{\omega}\cdot\nabla\boldsymbol{U}$引起的.

\textbf{(2)体积膨胀项$-\boldsymbol{\omega}\,(\nabla\cdot\boldsymbol{U})$}

由于$\nabla\cdot\boldsymbol{U}$是流体微团的体积膨胀率，对于一个在有势力场中运动的理想流体微团，它的角动量应是守恒的，流体微团体积膨胀导致它的转动惯量的增加，从而旋转角速度减小，也即涡量的减小.反之，流体微团收缩（$\nabla\cdot\boldsymbol{U}<0$）将导致涡量的增加.

\textbf{(3)有旋质量力场$\nabla\times\boldsymbol{f}$}

有旋的质量力场（例如旋转参照系统中的科氏力）能使流体微团发生旋转，也即产生涡量.如果质量力场是有势的，$\boldsymbol{f}=-\nabla\Pi$，这时质量力场对涡量的生成就没有贡献，因为$\nabla\times\boldsymbol{f}=\nabla\times(\nabla\Pi)=\boldsymbol{0}$.后面将举例说明非有势力场产生旋涡.

\textbf{(4)斜压流场的涡量生成项$\nabla\rho\times\nabla p/\rho^{2}$}

在斜压流场中$\nabla\rho\times\nabla p/\rho^{2}\neq\boldsymbol{0}$.例如有密度分层的大气、曲激波后的高速气流中都存在因斜压引起的涡量生成.对于正压流场，$\rho=\rho\,(p)$，因而$\nabla\rho\times\nabla p/\rho^{2}=\boldsymbol{0}$，这时压强梯度项$-\nabla p/\rho$也是有势的，所以它对涡量的生成没有贡献.

\subsection*{3. 斜压和科氏力产生有旋流场的实例：贸易风}

在地球东半球赤道附近，总是有一股从东北到西南的气流，在气象学上称为贸易风.它可以用理想流体的涡动力学理论来解释.

为了分析简明起见，假定地球是圆的，大气是完全气体.因此重力场中静止大气的等压面是同心圆球.由于赤道和北极区的太阳辐射强度不同，赤道附近的气温比北极区气温高，根据完全气体的状态方程，在相同的压强下，温度越高密度越低.于是，大气的密度从北极向赤道方向逐渐减小，造成等压面和等温面不重合，如图4.14所示.另一方面，大气的压强和密度都随高度而下降，因此压强梯度指向地心，而密度梯度指向地心偏西，形成斜压状态.根据压强梯度和密度梯度的方向，可以确定$\nabla\rho\times\nabla p/\rho^{2}$指向经度方向.根据涡动力学理论（式(4.68)），涡量在经度方向，随着时间的发展，在大圆上产生一环流，如图4.14所示，在高空该气流由南往北，在极地气流由上而下，在地面气流由北向南，在赤道气流由下而上.

进一步考察非有势力场中涡的产生.由于地球的自转，当大气相对于地球由北向南运动时，气体受到离心惯性力和科氏力的作用.科氏力是非有势的，并且
\[
\boldsymbol{f}_r=-2\,\boldsymbol{\omega}\times\boldsymbol{U}
\]
式中$\boldsymbol{\omega}=\omega\boldsymbol{k}$是地球的自转角速度向量，$\boldsymbol{k}$是地轴方向，由南极指向北极；$\boldsymbol{U}$是由斜压大气产生的由北向南气流，即$\boldsymbol{U}=U\boldsymbol{e}_\theta$（参见图4.15）.于是，科氏力
\[
\boldsymbol{f}_r=-2\,\boldsymbol{\omega}\times\boldsymbol{U}=-2\,\omega U\,\boldsymbol{k}\times\boldsymbol{e}_\theta=-2\,\omega U\,\boldsymbol{e}_\phi
\]
式中$\boldsymbol{e}_\phi$是纬度线由西向东的单位向量.于是从质点动力学的角度，不难理解科氏力将导致在地球表面由东向西的气流.

综合以上大气斜压和地球自转的联合作用，在地球赤道附近会有一股由东北到西南的气流，这就是贸易风.

\section*{习题}\addcontentsline{toc}{section}{习题}

\noindent\textbf{4.1}\quad 写出图4.16所示瞬时内圆柱和外圆柱壳上的理想流体的物面条件，坐标取在不动的外圆柱壳上.

\noindent\textbf{4.2}\quad 物体在无界理想流体中运动，分别采用下列坐标系建立物面条件，如图4.17所示：

（1）在固定于地面的直角坐标系$Oxyz$中讨论绝对运动；

（2）在固定于物体上的直角坐标系$O'x'y'z'$中讨论相对运动.

\noindent\textbf{4.3}\quad 证明不可压理想流体平面流动质量力有势时有：$\partial\omega^{2}/\partial t+\nabla\cdot(\omega^{2}\boldsymbol{U})=0$，$\omega$是平面流体运动的旋度.

\noindent\textbf{4.4}\quad 设有无穷远处速度分布为$U_\infty=ky$，$V_\infty=0$，$W_\infty=0$的理想均质不可压流体，绕过固定不动半径为$R$的无限长圆柱体，若质量力有势，试求涡量场.

\noindent\textbf{4.5}\quad 如图4.18所示，欲使离心水泵给图示管路系统提供$28\times10^{-3}\,\mathrm{m^{3}/s}$水量，需要多大理论水头（理论水头$=U^{2}/2g+H$）？

\noindent\textbf{4.6}\quad 如图4.19所示，半径为$a$的球体一边以$a(t)$膨胀，一边又以$U(t)$沿$Oz$负轴方向在无界原静止的不可压缩理想流体中作直线运动.求球面上的压力分布及所受的合力.设无穷远处静止流体的压力为$p_\infty$，密度为$\rho$.

\noindent\textbf{4.7}\quad 图4.20所示半径为$a$的二维圆柱，在无界原静止不可压理想流体中围绕$O$点作圆周运动，角速度为$\omega(t)$，旋转半径为$L$，证明物面压力分布为
\[
\frac{p}{\rho}=\frac{p_\infty}{\rho}+\dot{\omega}aL\sin\varepsilon+\frac{\omega^{2}L^{2}}{2}\,(1-4\cos^{2}\varepsilon)-\omega^{2}La\cos\varepsilon
\]
式中$\rho$为流体密度，$p_\infty$为无穷远处压力，已知$\Gamma=0$，忽略质量力.

\noindent\textbf{4.8}\quad 如图4.21所示，有一往复式活塞油泵，曲柄旋转角速度$\omega=30\,\mathrm{rad/s}$，缸内油的流动近似虚线所示一维流动.从1-1截面到2-2截面面积$A$的变化近似为$A=(5-40x)\times10^{-3}\,(\mathrm{m^{2}})$.已知活塞速度$U=0.5\omega(\sin\theta+0.1\sin 2\theta)$，$\theta=\omega t$.设油的比重为0.9，试求当$\theta=45^{\circ}$时，

（1）$a,b,c$三点的流体加速度.

（2）1-1截面与2-2截面上的油压差.

\noindent\textbf{4.9}\quad 敞口矩形水箱水面面积$F=0.5\,\mathrm{m^{2}}$，底部开一个直径为$10\,\mathrm{mm}$的小孔向大气泄水，流量系数为0.6，若初始水深为$1\,\mathrm{m}$，求泄空时间.又若水深保持不变，流动是定常的，泄出同样水量所需时间又为多少？

\noindent\textbf{4.10}\quad 图4.22所示等截面V形管内盛液体，液柱长度为$L$，当V形管作微小晃动后，求液面自由振荡的周期.

\noindent\textbf{4.11}\quad 在充满整个空间的原静止的不可压流体中，有一半径为$a$的空心球体突然破碎.设无穷远处流体压力为$p_\infty$，密度为$\rho$，质量力不计，求流体充满整个球体体积所需时间的积分表达式.

\noindent\textbf{4.12}\quad 原静止无界理想均质不可压缩流体中，有一球形物体，其半径以$R_b=1+3t$的规律膨胀（$t$为时间）（图4.23），求压力场.已知无穷远处压力为$p_\infty$，密度为$\rho$，且假定质量力可忽略.

\section*{思考题}\addcontentsline{toc}{section}{思考题}

\noindent\textbf{S4.1}\quad 圆球在牛顿型流体中突然移动，在启动瞬间绕圆球的流线谱是否和理想流体绕圆球的流线谱相同？为什么？如果没有任何质量力作用时，在理想流体中圆球等速膨胀或等速收缩时，球面上的压强如何变化？

\noindent\textbf{S4.2}\quad 在圆柱型桶中充满牛顿型流体，当圆桶等角速度旋转时，桶内的流体是否和圆桶一起转动？这时桶内的流动是否为有旋运动？这与Kelvin定理有矛盾吗？

\noindent\textbf{S4.3}\quad 为什么可以用“瞬时绝对坐标系”在随物体运动的坐标系中研究该瞬间流体的绝对运动？

\noindent\textbf{S4.4}\quad 计算在海浪作用下圆柱形电缆运动时是否需要计入附加质量？